\documentclass[11pt,a4paper]{article}
\usepackage[T1]{fontenc}
\usepackage{lmodern,amsmath,amssymb,amsthm,mathtools}
\usepackage[margin=26mm]{geometry}
\usepackage{microtype,booktabs,array,xcolor,listings,hyperref}
\newcolumntype{P}[1]{>{\raggedright\arraybackslash}p{#1}}
\definecolor{navy}{RGB}{25,48,69}
\hypersetup{colorlinks=true,linkcolor=navy,urlcolor=navy,citecolor=navy,
  pdftitle={Parameter Monotonicity of Mittag-Leffler Functions: Certified Bounds and Threshold Geometry},
  pdfauthor={Prepared for Keio}}
\theoremstyle{plain}
\newtheorem{theorem}{Theorem}[section]
\newtheorem{lemma}[theorem]{Lemma}
\newtheorem{corollary}[theorem]{Corollary}
\theoremstyle{remark}
\newtheorem*{remark}{Remark}
\newcommand{\E}{\mathbb E}
\newcommand{\R}{\mathbb R}
\newcommand{\Z}{\mathbb Z}
\newcommand{\C}{\mathbb C}
\DeclareMathOperator{\Log}{Log}
\DeclareMathOperator{\Res}{Res}
\lstset{language=Python,basicstyle=\ttfamily\scriptsize,columns=fullflexible,keepspaces=true,
  breaklines=true,showstringspaces=false,frame=single,framerule=0.3pt,rulecolor=\color{gray},
  upquote=true,tabsize=4,xleftmargin=2pt,xrightmargin=2pt}
\setlength{\emergencystretch}{2em}
\allowdisplaybreaks

\title{\color{navy}Parameter Monotonicity of Mittag--Leffler Functions:\\
Certified Bounds, Threshold Geometry,\\ and Fractional Comparison}
\author{Prepared for Keio\thanks{Revised from the supplied unified manuscript of
15 November 2026, retaining its two proofs and reproducible certificates.
The review, stronger certificates, and extension theorems were prepared in this
revision. No external refereeing or historical-priority certification is claimed.}}
\date{15 November 2026\\\small Reviewed and extended computer-assisted research manuscript}

\begin{document}
\maketitle

\begin{abstract}
We review and extend a computer-assisted proof of parameter monotonicity for
$U_\alpha(x)=E_\alpha(x^\alpha)$. Two reproduced contour arguments establish
the expectation of Ferreira and Simon on $(2,4)$. Strengthened Arb certificates give
$8.803377\le\alpha^*<8.8033777$, where $\alpha^*$ is the maximal initial
monotonicity endpoint. We also prove that for every $\alpha\ge8.8033777$ there is
an $x>0$ with $\partial_\alpha U_\alpha(x)>0$, including the entire unbounded
range $\alpha\ge9$ by a first-term dominance argument. For the two-parameter
function, the maximal endpoint $A(\beta)$ is finite, strictly increasing and
right-continuous for $\beta\ge1$, and is attained at an interior tangency.
An additional certificate proves monotonicity through order $9$ for every
$\beta\ge2$. We transfer these results to complete monotonicity of explicit
Laplace kernels and comparison for fractional Volterra equations.
All computer-assisted claims include reproducible code. The numerical value
$8.80337769817047\ldots$ is a candidate, not an exact certified threshold;
uniqueness and nondegeneracy of its contact remain unproved.
\end{abstract}

\tableofcontents

\section{Introduction}\label{sec:intro}
\subsection{Background}
For $\alpha,\beta>0$ let
\[
 E_{\alpha,\beta}(z)=\sum_{n=0}^{\infty}\frac{z^n}{\Gamma(\beta+\alpha n)},\qquad
 E_\alpha=E_{\alpha,1},\qquad U_\alpha(x)=E_\alpha(x^\alpha)\quad(x>0).
\]
The scaling $x^\alpha$ puts all orders on one exponential scale: $U_\alpha(x)\sim e^x/\alpha$
as $x\to\infty$, and
\[
 U_1(x)=e^x,\quad U_2(x)=\cosh x,\quad
 U_3(x)=\tfrac13\bigl(e^x+2e^{-x/2}\cos\tfrac{\sqrt3\,x}{2}\bigr),\quad
 U_4(x)=\tfrac12(\cosh x+\cos x).
\]
Equivalently, $U_\alpha$ solves the Caputo problem $D^\alpha u=u$, $u(0)=1$,
$u^{(j)}(0)=0$ for $1\le j<\lceil\alpha\rceil$ \cite{GKMR}. Monotonicity in $\alpha$ therefore
compares solutions of fractional growth equations of different orders.

Ferreira and Simon \cite[Theorem~A]{FS} proved that $\alpha\mapsto U_\alpha(x)$ decreases on
$(0,2)$ for every $x>0$, through convex ordering of an associated family of random variables.
In \cite[Remark~1(a)]{FS} they discuss larger orders. They expect, supported by simulations, that
the decrease persists on $(2,4)$, and they check $U_4<U_3<U_2$ from the closed forms above. They
also explain why their probabilistic argument does not reach $\alpha>2$, and observe numerically
that $U_{25}(30)<U_{26}(30)$, so monotonicity fails for some large $\alpha$. Two questions remain:
the conjecture on $(2,4)$, and the exact range of $\alpha$ on which the decrease holds for all $x>0$.

\subsection{Main results}
Write $H_\alpha(x)=-\partial_\alpha U_\alpha(x)$ and $\psi=\Gamma'/\Gamma$.

\begin{theorem}\label{thm:24}
For every $x>0$ and every $\alpha\in[2,4]$, $H_\alpha(x)>0$. In particular
$\alpha\mapsto E_\alpha(x^\alpha)$ is strictly decreasing on $[2,4]$, and the conjecture of
\cite[Remark~1(a)]{FS} holds.
\end{theorem}

\begin{theorem}\label{thm:sharp}
\begin{enumerate}
 \item[(i)] For every $x>0$, the map $\alpha\mapsto E_\alpha(x^\alpha)$ is strictly decreasing on
 $(0,8.803377]$.
 \item[(ii)] $\partial_\alpha E_\alpha(x^\alpha)>0$ at $(\alpha,x)=(8.8033777,\,10.7538996)$.
\end{enumerate}
Consequently the threshold
\[
 \alpha^*:=\sup\bigl\{A>0:\ \alpha\mapsto E_\alpha(x^\alpha)\ \text{is non-increasing on }(0,A]
 \text{ for every }x>0\bigr\}
\]
satisfies $8.803377\le\alpha^*<8.8033777$. The candidate $8.80337769817\ldots$ is numerical only
(Section~\ref{subsec:critical}); exact identification and uniqueness are not asserted.
\end{theorem}

\begin{corollary}\label{cor:beta}
For every $\beta\ge1$ and $x>0$, the map $\alpha\mapsto E_{\alpha,\beta}(x^\alpha)$ is strictly
decreasing on $(0,8.803377]$.
\end{corollary}

\begin{corollary}\label{cor:bounds}
For every $x>0$,
\begin{gather*}
 \frac{\cosh x+\cos x}{2}<E_\alpha(x^\alpha)<\cosh x\qquad(2<\alpha<4),\\
 E_8(x^8)<E_\alpha(x^\alpha)<\frac{\cosh x+\cos x}{2}\qquad(4<\alpha<8),
\end{gather*}
where
$E_8(x^8)=\tfrac14\bigl(\cosh x+\cos x+2\cos\tfrac{x}{\sqrt2}\cosh\tfrac{x}{\sqrt2}\bigr)$.
\end{corollary}

\subsection{Structure of the proofs}
Every argument below uses the Laplace transform of $H_\alpha$,
\[
 K_\alpha(s)=\frac{s^{\alpha-1}\Log s}{(s^\alpha-1)^2}\qquad(\Re s>1),
\]
and its Bromwich inversion (Lemmas~\ref{lem:laplace} and~\ref{lem:inversion}); the proofs differ
in where the contour is moved.
\begin{itemize}
 \item Proof~A (Section~\ref{sec:A}) shifts the vertical line to $\Re s=1/2$. This crosses only
 the pole at $s=1$, but the remainder is of size $e^{x/2}$, so a compact rectangle in $(\alpha,x)$
 is certified.
 \item Proof~B (Section~\ref{sec:B}) rotates the contour onto two rays in the left half-plane and
 keeps the residues at $e^{\pm2\pi i/\alpha}$ explicitly. The remainder then stays bounded in $x$,
 and the certified part is one-dimensional.
 \item For $\alpha>4$ more residues enter (Section~\ref{sec:sharp}). The rays of Proof~B are then
 combined with a two-dimensional certificate in the style of Proof~A, which is needed only near the
 critical point.
\end{itemize}
Small $x$ is handled in every case by termwise positivity of the series for $H_\alpha$
(Lemma~\ref{lem:small}). Table~\ref{tab:cert} lists the finite certified computations.

\begin{table}[ht]\centering\small\setlength{\tabcolsep}{4pt}
\begin{tabular}{P{2.7cm}P{5.0cm}P{3.4cm}P{3.2cm}}
\toprule
Step & Region & Certificate & Certified margin\\
\midrule
Lemma~\ref{lem:certA} (Proof A) & $[2,4]\times[\frac52,16]$ & 2-D, 16{,}749 rectangles &
 $e^{-x}H_{\alpha,24}>10^{-5}$\\
Lemma~\ref{lem:certB} (Proof B) & $\alpha\in[2,4]$ & 1-D, 2{,}000 intervals & $G\le0.1999$\\
Lemma~\ref{lem:tail16} & $\alpha\in[4,8.803377]$, $x\ge16$ & 1-D, 960 intervals & $G\le0.501$\\
Lemma~\ref{lem:middle} & $\alpha\in[4,8.803377]$, $e^{\psi(1+\alpha)}\le x\le16$ & 2-D, 413 boxes &
 $Q\ge6.12\cdot10^{-9}$\\
Lemma~\ref{lem:witness} & $(\alpha,x)=(8.8033777,10.7538996)$ & single point &
 $\partial_\alpha U_\alpha\ge3.2560\cdot10^{-7}$\\
\bottomrule
\end{tabular}
\caption{Finite certified computations. All use Arb ball arithmetic \cite{Arb} through
python-flint~0.9.0 \cite{flint}; the programs are listed in Appendix~\ref{app:code}. Here
$G<1$ and $Q>0$ are the sufficient conditions of Lemma~\ref{lem:bounds}.}\label{tab:cert}
\end{table}

\subsection{Scope and review findings}
The supplied draft's main proof survives mathematical review and successful
replay of all three certificate programs. This revision repairs the proof of
the infinite witness tail and replaces a rounded-down bound $0.5005$ by the
valid bound $0.501$. It also removes an unsupported identification of a
numerical double root with the exact first threshold. Beyond the original
result, Sections~\ref{sec:no-recovery}--\ref{sec:transfers} prove global
failure at all sufficiently large orders, a strict hierarchy in the second
parameter, and applications. The source expectation on $(2,4)$ is documented
in \cite[Remark 1(a)]{FS}; the new extension theorems are not presented as
solutions of additional independently documented historical conjectures.
This is a reproducible research draft, not a proof-assistant formalization or
an externally refereed result.

\section{Analytic preliminaries}\label{sec:prelim}
Throughout, $\Log s$ and $s^\beta=e^{\beta\Log s}$ are principal branches on
$\C\setminus(-\infty,0]$.

\begin{lemma}[derivative series; small $x$]\label{lem:small}
For $x>0$ and $\alpha>0$,
\begin{equation}\label{eq:series}
 H_\alpha(x)=\sum_{n=1}^{\infty}\frac{n\,x^{\alpha n}}{\Gamma(1+\alpha n)}
 \bigl(\psi(1+\alpha n)-\log x\bigr),
\end{equation}
with locally uniform convergence in $(\alpha,x)$. If $0<x\le e^{\psi(1+\alpha)}$, every term is
nonnegative and every term with $n\ge2$ is positive; hence $H_\alpha(x)>0$.
\end{lemma}
\begin{proof}
The series for $U_\alpha$ and its termwise $\alpha$-derivative converge locally uniformly because
$\Gamma(1+\alpha n)$ grows faster than any exponential in $n$, and
$\partial_\alpha\bigl[x^{\alpha n}/\Gamma(1+\alpha n)\bigr]
 =n\,x^{\alpha n}\bigl(\log x-\psi(1+\alpha n)\bigr)/\Gamma(1+\alpha n)$.
Since $\psi'(y)=\sum_{j\ge0}(y+j)^{-2}>0$, we have
$\psi(1+\alpha n)\ge\psi(1+\alpha)\ge\log x$ for $n\ge1$, with strict first inequality for $n\ge2$.
\end{proof}

For $\alpha\in[2,8.803377]$ the threshold $e^{\psi(1+\alpha)}$ increases from
$e^{\psi(3)}=2.5162\ldots$ to $e^{\psi(9.8033)}=9.3077\ldots$; at $\alpha=4$ it equals
$e^{\psi(5)}=4.5091\ldots$.

\begin{lemma}[Laplace transform]\label{lem:laplace}
For $\alpha\ge2$ and $\Re s>1$, $\displaystyle\int_0^\infty e^{-sx}H_\alpha(x)\,dx=K_\alpha(s)$.
\end{lemma}
\begin{proof}
Let $c=\Re s>1$, $n\ge1$ and $y=1+\alpha n\ge1+2n$. If $Y$ has the gamma law of shape $y$ and unit
rate, the substitution $x=u/c$ gives
\begin{align*}
 \int_0^\infty e^{-cx}\frac{n\,x^{\alpha n}}{\Gamma(1+\alpha n)}\bigl|\psi(y)-\log x\bigr|\,dx
 &=n\,c^{-y}\,\E\bigl|\psi(y)-\log Y+\log c\bigr|\\
 &\le n\,c^{-1-2n}\bigl(\psi(y)+y+(y-1)^{-1}+\log c\bigr),
\end{align*}
using $0<\psi(y)=\E\log Y$ and $\E|\log Y|\le\E(Y+Y^{-1})=y+(y-1)^{-1}$. The bound is
$O(n^2c^{-2n})$ locally uniformly in $\alpha$, hence summable, and termwise integration is
justified. Differentiating $\int_0^\infty e^{-sx}x^\nu\,dx=\Gamma(\nu+1)s^{-\nu-1}$ in $\nu$ gives
$\int_0^\infty e^{-sx}x^\nu\bigl(\psi(\nu+1)-\log x\bigr)dx/\Gamma(\nu+1)=s^{-\nu-1}\Log s$.
With $\nu=\alpha n$ and $|s^{-\alpha}|<1$,
$\sum_{n\ge1}n\,s^{-1-\alpha n}\Log s=s^{\alpha-1}\Log s/(s^\alpha-1)^2$.
\end{proof}

\begin{lemma}[inversion]\label{lem:inversion}
For $\alpha\ge2$, $c>1$ and $x>0$,
\[
 H_\alpha(x)=\frac{1}{2\pi i}\int_{c-i\infty}^{c+i\infty}e^{xs}K_\alpha(s)\,ds,
\]
and the integral converges absolutely.
\end{lemma}
\begin{proof}
On $\Re s=c$ we have $|K_\alpha(c+it)|=O(|t|^{-\alpha-1}\log|t|)$, so $t\mapsto K_\alpha(c+it)$ is
integrable. By Lemma~\ref{lem:laplace} it is the Fourier transform of
$f(x)=e^{-cx}H_\alpha(x)\mathbf 1_{x>0}\in L^1(\R)$. Fourier inversion recovers $f$ almost
everywhere; both sides are continuous on $(0,\infty)$.
\end{proof}

\begin{lemma}[poles and residues]\label{lem:poles}
On $\C\setminus(-\infty,0]$, $K_\alpha$ is meromorphic, with poles exactly at $w_k=e^{iu_k}$,
$u_k=2\pi k/\alpha$, $k\in\Z$, $|u_k|<\pi$. The pole $w_0=1$ is simple, the others are double, and
\[
 \Res_{s=w_k}e^{xs}K_\alpha(s)=\frac{e^{xw_k}}{\alpha^2}\bigl(1+iu_k\,x\,w_k\bigr).
\]
\end{lemma}
\begin{proof}
$s^\alpha=1$ if and only if $\alpha\Log s\in2\pi i\Z$, and these zeros are simple because
$(s^\alpha)'=\alpha s^{\alpha-1}\ne0$. Put $g(s)=e^{xs}s^{\alpha-1}\Log s$ and $F(s)=s^\alpha-1$.
At a simple zero $p$ of $F$, $\Res(g/F^2)=g'(p)/F'(p)^2-g(p)F''(p)/F'(p)^3$. With $p^\alpha=1$,
$F'(p)=\alpha/p$, $F''(p)=\alpha(\alpha-1)/p^2$ and $\Log p=iu_k$, this equals
$e^{xp}(1+x\,p\Log p)/\alpha^2$. At $p=1$ the zero of $\Log s$ cancels one order and the same
formula gives $e^x/\alpha^2$.
\end{proof}

\begin{lemma}[rotated rays]\label{lem:rays}
Let $\alpha\ge2$, $x>0$ and $\varphi\in(\pi/2,\pi)$ with $\alpha\varphi\notin2\pi\Z$, and put
$K_\varphi=\#\{k\ge1:u_k<\varphi\}$. Then
\begin{equation}\label{eq:rays}
 H_\alpha(x)=\frac{e^x+S_\varphi(\alpha,x)}{\alpha^2}+J_\varphi(\alpha,x),
\end{equation}
where
\[
 S_\varphi(\alpha,x)=\sum_{k=1}^{K_\varphi}2e^{x\cos u_k}
 \Bigl[\cos(x\sin u_k)-u_kx\,\sin(x\sin u_k+u_k)\Bigr],
\]
\[
 J_\varphi(\alpha,x)=\frac1\pi\,\Im\int_0^\infty e^{xre^{i\varphi}}\,r^{\alpha-1}e^{i\alpha\varphi}\,
 \frac{\log r+i\varphi}{\bigl(r^\alpha e^{i\alpha\varphi}-1\bigr)^2}\,dr .
\]
\end{lemma}
\begin{proof}
Let $\Gamma_\varphi$ run from $\infty e^{-i\varphi}$ to $0$ and on to $\infty e^{i\varphi}$, and
consider the region $\{|\arg s|<\varphi,\ \Re s<c\}$ between $\Gamma_\varphi$ and the line of
Lemma~\ref{lem:inversion}; the whole line lies in the sector because $|\arg(c+it)|<\pi/2$.
$K_\alpha$ has no pole on $\Gamma_\varphi$, since $s^\alpha=1$ on $\arg s=\pm\varphi$ would force
$\alpha\varphi\in2\pi\Z$. On arcs $|s|=R$ with $\Re s\le c$ we have $|e^{xs}|\le e^{cx}$ and
$|K_\alpha(s)|=O(R^{-\alpha-1}\log R)$, so these arcs contribute nothing as $R\to\infty$. Near
$s=0$, $|K_\alpha(s)|=O(|s|^{\alpha-1}|\log|s||)$, so a small arc around the origin also
contributes nothing. Every pole has modulus $1<c$, so the residue theorem gives
\[
 H_\alpha(x)=\sum_{|u_k|<\varphi}\Res_{s=w_k}e^{xs}K_\alpha(s)
 +\frac1{2\pi i}\int_{\Gamma_\varphi}e^{xs}K_\alpha(s)\,ds .
\]
By Lemma~\ref{lem:poles} the residues at $w_{\pm k}$ are complex conjugate, and their sum is the
$k$-th term of $S_\varphi/\alpha^2$. Since $K_\alpha(\bar s)=\overline{K_\alpha(s)}$, the ray
integral equals $J_\varphi$. The argument stays inside the open sector $|\arg s|<\varphi<\pi$ and
never meets the branch cut, so it holds for even integer $\alpha$ as well.
\end{proof}

\begin{lemma}[explicit bounds]\label{lem:bounds}
In the setting of Lemma~\ref{lem:rays} put $\lambda=-x\cos\varphi>0$ and $\gamma=\alpha\varphi$,
and let $m^2=1$ if $\cos\gamma\le0$ and $m^2=\sin^2\gamma$ otherwise. Then
$|J_\varphi(\alpha,x)|\le B_\varphi(\alpha,x)$, where
\[
 B_\varphi(\alpha,x)=\frac{1}{\pi m^2}\Bigl[\frac{\Gamma(\alpha)}{\lambda^\alpha}
 \bigl(\varphi+(\psi(\alpha)-\log\lambda)_+\bigr)
 +2\min\Bigl(\frac1{\alpha^2},\frac{\Gamma(\alpha-1)}{e\,\lambda^{\alpha-1}}\Bigr)\Bigr],
\]
and $|e^{-x}S_\varphi(\alpha,x)|\le\sum_{k=1}^{K_\varphi}2e^{-x(1-\cos u_k)}(1+u_kx)$.
Consequently
\[
 \alpha^2e^{-x}H_\alpha(x)\ \ge\ Q_\varphi(\alpha,x):=1+e^{-x}S_\varphi(\alpha,x)
 -\alpha^2e^{-x}B_\varphi(\alpha,x)\ \ge\ 1-G_\varphi(\alpha,x),
\]
\[
 G_\varphi(\alpha,x):=\sum_{k=1}^{K_\varphi}2e^{-x(1-\cos u_k)}(1+u_kx)
 +\alpha^2e^{-x}B_\varphi(\alpha,x).
\]
In particular $H_\alpha(x)>0$ whenever $Q_\varphi(\alpha,x)>0$, and in particular whenever
$G_\varphi(\alpha,x)<1$.
\end{lemma}
\begin{proof}
On $s=re^{\pm i\varphi}$ we have $|e^{xs}|=e^{-\lambda r}$, $|s^{\alpha-1}|=r^{\alpha-1}$,
$|\log r\pm i\varphi|\le|\log r|+\varphi$ and
$|r^\alpha e^{\pm i\gamma}-1|^2=(r^\alpha-\cos\gamma)^2+\sin^2\gamma\ge m^2$. Hence
$|J_\varphi|\le(\pi m^2)^{-1}\int_0^\infty e^{-\lambda r}r^{\alpha-1}(|\log r|+\varphi)\,dr$. Write
$|\log r|=\log r+2\log(1/r)\mathbf 1_{r<1}$ and use
$\int_0^\infty e^{-\lambda r}r^{\alpha-1}\log r\,dr=\Gamma(\alpha)\lambda^{-\alpha}(\psi(\alpha)-\log\lambda)$.
The remaining integral $L=\int_0^1e^{-\lambda r}r^{\alpha-1}\log(1/r)\,dr$ is at most
$\int_0^1r^{\alpha-1}\log(1/r)\,dr=\alpha^{-2}$, and, since $\log(1/r)\le 1/(er)$, at most
$\Gamma(\alpha-1)/(e\lambda^{\alpha-1})$. Finally
$|e^{xw}(1+iuxw)|\le e^{x\cos u}(1+ux)$ for $w=e^{iu}$.
\end{proof}

\begin{lemma}[monotonicity in $x$]\label{lem:mono}
Let $2\le\alpha\le8.8033777$ and let $\varphi$ be as in Lemma~\ref{lem:rays}. Then
$x\mapsto G_\varphi(\alpha,x)$ is strictly decreasing on $[4.1,\infty)$, and on $[1,\infty)$ if
$\alpha\le4$.
\end{lemma}
\begin{proof}
For $c=\cos u$, the derivative of $e^{-x(1-c)}(1+ux)$ is $e^{-x(1-c)}\bigl(u-(1-c)(1+ux)\bigr)$,
which is negative for $x>(1-c)^{-1}-u^{-1}$. Here $u\ge u_1=2\pi/\alpha\ge0.7137$, so
$1-c\ge1-\cos0.7137>0.244$ and $(1-c)^{-1}<4.1$; if $\alpha\le4$ then $c\le0$ and $(1-c)^{-1}\le1$.
In $B_\varphi$, $\lambda$ is proportional to $x$, and the functions $\lambda^{-\alpha}$,
$\lambda^{-\alpha}(\psi(\alpha)-\log\lambda)_+$ (derivative
$-\lambda^{-\alpha-1}[1+\alpha(\psi(\alpha)-\log\lambda)]<0$ where it is positive) and the minimum are
nonincreasing in $\lambda$. The factor $e^{-x}$ makes $e^{-x}B_\varphi$ strictly decreasing.
\end{proof}

\section{Proof A of Theorem~\ref{thm:24}: the line \texorpdfstring{$\Re s=1/2$}{Re s = 1/2}}\label{sec:A}
This section follows the draft of 14 November 2026. Its compact certificate was re-implemented
independently (Appendix~\ref{app:A}) and reproduces the draft's data exactly.

\begin{lemma}\label{lem:line}
Uniformly for $\alpha\in[2,4]$ and $x>0$, $\bigl|H_\alpha(x)-e^x/\alpha^2\bigr|<100\,e^{x/2}$.
In particular $H_\alpha(x)>0$ for $x\ge16$.
\end{lemma}
\begin{proof}
For $\alpha\in[2,4]$ the only pole of $K_\alpha$ in $\Re s>0$ is $s=1$: there
$|\arg s|<\pi/2$, so $\alpha\arg s$ cannot be a nonzero multiple of $2\pi$. At $\alpha=4$ the
further roots $\pm i$ lie to the left of both lines below. Moving the line of
Lemma~\ref{lem:inversion} to $\Re s=1/2$ crosses only $s=1$, and the horizontal segments vanish by
the decay of $K_\alpha$. Hence
\[
 H_\alpha(x)=\frac{e^x}{\alpha^2}+\frac1{2\pi}\int_{-\infty}^{\infty}
 e^{(1/2+it)x}K_\alpha\bigl(\tfrac12+it\bigr)\,dt .
\]
Let $s=\frac12+it$, $r=|s|$, and take $t\ge0$ (conjugation handles $t<0$).
\begin{enumerate}
 \item If $0\le t\le\frac12$: $r\le2^{-1/2}$, $|s^\alpha-1|\ge1-r^\alpha\ge\frac12$,
 $|\Log s|\le\log2+\frac\pi4<\frac32$ and $r^{\alpha-1}\le1$, so $|K_\alpha(s)|<6$.
 \item If $\frac12\le t\le2$: $\frac\pi4\le\arg s\le\arctan4$, and
 $\frac\pi2\le\alpha\arg s<\frac{11\pi}6$ because $\arctan\frac14\ge\frac4{17}>\frac\pi{24}$. If
 $\cos(\alpha\arg s)\le0$ the distance from $s^\alpha$ to $1$ is at least $1$; otherwise
 $\alpha\arg s\in(\frac{3\pi}2,\frac{11\pi}6)$ and $|\sin(\alpha\arg s)|\ge\frac12$. Thus
 $|s^\alpha-1|\ge\frac12$. With $r\le\sqrt{17}/2$, $r^{\alpha-1}\le\max(r,r^3)<9$ and
 $|\Log s|<3$ we get $|K_\alpha(s)|<108<144$.
 \item If $t\ge2$: $r^\alpha\ge4$, $|s^\alpha-1|\ge\frac34r^\alpha$, $\log(r/t)\le\frac1{8t^2}<\frac14$
 and $|\arg s|<\frac\pi2<2$, so $|K_\alpha(s)|<2t^{-3}(\log t+\frac94)$.
\end{enumerate}
Therefore
$\int_\R|K_\alpha(\frac12+it)|\,dt<2\bigl(3+216+\frac{4\log2+11}{16}\bigr)<440$, and
$|H_\alpha(x)-e^x/\alpha^2|<\frac{440}{2\pi}e^{x/2}<100e^{x/2}$. For $x\ge16$,
$H_\alpha(x)>e^x/16-100e^{x/2}>0$ because $e^8>1600$.
\end{proof}

\begin{lemma}\label{lem:scalar}
$\psi(3)>\log\frac52$ and $\psi(51)>\log16$. Consequently, for $\alpha\in[2,4]$, every term
of~\eqref{eq:series} is positive when $0<x\le\frac52$, and every term with $n\ge25$ is positive
when $x\le16$. In particular $H_\alpha(x)\ge H_{\alpha,24}(x)$ for $x\in[\frac52,16]$, where
\[
 H_{\alpha,24}(x)=\sum_{n=1}^{24}\frac{n\,x^{\alpha n}}{\Gamma(1+\alpha n)}
 \bigl(\psi(1+\alpha n)-\log x\bigr).
\]
\end{lemma}
\begin{proof}
Arb gives $\psi(3)-\log\frac52\in[0.00649360\pm3.3\cdot10^{-9}]$ and
$\psi(51)-\log16\in[1.14940\pm10^{-6}]$ (Appendix~\ref{app:A}); note $\psi(3)=\frac32-\gamma_E$.
The rest follows from $1+\alpha n\ge3$, from $1+\alpha n\ge51$ for $n\ge25$, and from the
monotonicity of $\psi$.
\end{proof}

\begin{lemma}[compact certificate]\label{lem:certA}
For every $(\alpha,x)\in[2,4]\times[\frac52,16]$, $e^{-x}H_{\alpha,24}(x)>10^{-5}$.
\end{lemma}
\begin{proof}
The program of Appendix~\ref{app:A} starts from the exact rational rectangle
$[2,4]\times[\frac52,16]$. On each rectangle it encloses
\[
 \sum_{n=1}^{24}n\exp\bigl(nA\log X-\log\Gamma(1+nA)-X\bigr)\bigl(\psi(1+nA)-\log X\bigr)
\]
in 96-bit Arb ball arithmetic, with $A$ and $X$ balls containing the parameter intervals, and
accepts the rectangle if the lower end of the ball exceeds the exact rational $10^{-5}$.
Otherwise it bisects $\alpha$ when eight times its width exceeds the $x$-width, and $x$ otherwise.
All endpoints are exact rationals, each split replaces a rectangle by two with the same union, and
a depth guard can only cause failure. The run terminates with an empty stack after
33{,}497 evaluations and accepts 16{,}749 rectangles of total exact area $27$, with maximum depth
$16$ (so $33497=2\cdot16749-1$). The smallest certified lower bound is $1.1489\cdot10^{-5}$.
\end{proof}

The draft reported the same rectangle data from its own program, together with a replay at 160
bits using direct powers and gamma values (34{,}246 subrectangles in 51{,}743 evaluations). That
replay is not part of this bundle.

\begin{proof}[Proof A of Theorem~\ref{thm:24}]
For $0<x\le\frac52$ use Lemma~\ref{lem:scalar}. For $x\in[\frac52,16]$, Lemmas~\ref{lem:scalar}
and~\ref{lem:certA} give $H_\alpha(x)\ge H_{\alpha,24}(x)>0$. For $x\ge16$ use
Lemma~\ref{lem:line}. These three ranges cover $x>0$ uniformly in $\alpha\in[2,4]$.
\end{proof}

\begin{remark}
The constant in Lemma~\ref{lem:line} is coarse. Numerically,
$\int_\R|K_\alpha(\frac12+it)|\,dt\le1.73$ for $\alpha\in[2,4]$, with the maximum at $\alpha=2$.
With this value and the $\alpha$-dependent threshold of Lemma~\ref{lem:small}, the ratio of the
remainder bound to $e^x/\alpha^2$ at $x=e^{\psi(1+\alpha)}$ is at most $0.31$. A certified version of
this one-dimensional integral bound would therefore remove the compact certificate from Proof~A.
This observation is not used.
\end{remark}

\section{Proof B of Theorem~\ref{thm:24}: rotated rays}\label{sec:B}
Use Lemma~\ref{lem:rays} with
\[
 \varphi=\tfrac{7\pi}{10}\ \text{on }[2,\tfrac52]\ (K_\varphi=0),\qquad
 \varphi=\tfrac{9\pi}{10}\ \text{on }[\tfrac52,4]\ (K_\varphi=1).
\]
On the first range $u_1=2\pi/\alpha\ge\frac{4\pi}5>\varphi$, so no residue besides $s=1$ enters. On
the second, $u_1\le\frac{4\pi}5<\varphi$, while $u_2=4\pi/\alpha\ge\pi$ is not a pole. In both
ranges $\alpha\varphi\notin2\pi\Z$, and $m^2\ge\frac12$ because $\gamma=\alpha\varphi$ lies in
$[1.4\pi,1.75\pi]$ and in $[2.25\pi,3.6\pi]$, respectively.

\begin{lemma}[one-dimensional certificate]\label{lem:certB}
Split $[2,\frac52]$ into $500$ and $[\frac52,4]$ into $1500$ equal intervals $[a,a']$. For each,
let $X$ be a certified lower bound of $e^{\psi(1+a)}$. Then
$\sup_{\alpha\in[a,a']}G_\varphi(\alpha,X)<1$. The largest certified upper bounds are $0.199825$
(on $[2.499,2.5]$) and $0.182407$ (on $[3.999,4]$).
\end{lemma}
\begin{proof}
The program of Appendix~\ref{app:B} evaluates $G_\varphi$ with $\alpha$ a ball containing $[a,a']$,
in 128-bit Arb arithmetic, with a certified lower bound of $m^2$. On every interval it also
asserts that $u_1$ lies on the stated side of $\varphi$ and that $X>1$.
\end{proof}

\begin{proof}[Proof B of Theorem~\ref{thm:24}]
Let $\alpha\in[2,4]$ and $x>0$. If $x\le e^{\psi(1+\alpha)}$, use Lemma~\ref{lem:small}. Otherwise,
for the interval $[a,a']\ni\alpha$ of Lemma~\ref{lem:certB}, monotonicity of $\psi$ gives
$x\ge e^{\psi(1+\alpha)}\ge X>2.5$. Lemmas~\ref{lem:mono} and~\ref{lem:certB} give
$G_\varphi(\alpha,x)\le G_\varphi(\alpha,X)<1$, and Lemma~\ref{lem:bounds} gives $H_\alpha(x)>0$.
\end{proof}

\begin{remark}
Proofs~A and~B share only Lemmas~\ref{lem:small}--\ref{lem:inversion}; their contour shifts,
bounds and certificates are independent. As a non-rigorous cross-check, formula~\eqref{eq:rays}
agrees with the series~\eqref{eq:series} to relative error at most $9.6\cdot10^{-41}$ in $36$ cases
($\alpha\in\{2.1,2.3,2.7,3,3.4,3.9\}$, $x\in\{0.7,2.5,6\}$, $\varphi\in\{0.75\pi,0.92\pi\}$,
40-digit arithmetic).
\end{remark}

\section{A certified transition window: proof of Theorem~\ref{thm:sharp}}\label{sec:sharp}
By \cite[Theorem~A]{FS} and Theorem~\ref{thm:24}, it remains to treat $\alpha\in[4,8.803377]$ and to
certify the witness of part~(ii).

\subsection{Ray angles}
For $\alpha>4$ the poles $w_{\pm1}$ move into the right half-plane, and further poles $w_{\pm2}$,
$w_{\pm3}$, $w_{\pm4}$ appear on the principal sheet as $\alpha$ passes $4$, $6$ and $8$. We fix one
ray angle on each of five ranges (Table~\ref{tab:rays}).

\begin{table}[ht]\centering\small
\begin{tabular}{lcccc}
\toprule
$\alpha$-range & $\varphi/\pi$ & $K_\varphi$ & last pole inside, $u_{K_\varphi}/\pi$ &
 first pole outside, $u_{K_\varphi+1}/\pi$\\
\midrule
$[4,5]$ & $0.65$ & 1 & $[0.40,0.50]$ & $[0.80,1.00]$\\
$[5,6]$ & $0.90$ & 2 & $[0.667,0.80]$ & $[1.00,1.20]$\\
$[6,7]$ & $0.76$ & 2 & $[0.571,0.667]$ & $[0.857,1.00]$\\
$[7,8]$ & $0.66$ & 2 & $[0.50,0.571]$ & $[0.75,0.857]$\\
$[8,8.803377]$ & $0.83$ & 3 & $[0.682,0.75]$ & $[0.909,1.00]$\\
\bottomrule
\end{tabular}
\caption{Ray angles for $\alpha\in[4,8.803377]$; $u_k=2\pi k/\alpha$. A pole with $u_k\ge\pi$ is not
on the principal sheet.}\label{tab:rays}
\end{table}

\begin{lemma}[geometry]\label{lem:geom}
On each range of Table~\ref{tab:rays} and for every $\alpha$ in it: $\varphi\in(\pi/2,\pi)$,
$u_k<\varphi$ for $1\le k\le K_\varphi$, $u_{K_\varphi+1}>\varphi$, and $m^2>\frac12$. In
particular $\alpha\varphi\notin2\pi\Z$, and Lemma~\ref{lem:rays} applies with the stated
$K_\varphi$.
\end{lemma}
\begin{proof}
The inequalities for $u_k$ are checked in ball arithmetic on each whole range, and $m^2>\frac12$
on $400$ equal sub-intervals of each range (Appendix~\ref{app:C}). If $\alpha\varphi=2\pi k$ then
$\varphi=u_k$, which the inequalities exclude.
\end{proof}

\begin{lemma}[large $x$]\label{lem:tail16}
For every $\alpha\in[4,8.803377]$, $G_\varphi(\alpha,16)\le0.501$. Hence $H_\alpha(x)>0$ for
$x\ge16$.
\end{lemma}
\begin{proof}
The bound is certified on $960$ equal sub-intervals: $200$ on each of the first four ranges and
$160$ on $[8,8.803377]$. The largest certified upper bound is $0.500531354737$ after upward rounding. The second statement
follows from Lemmas~\ref{lem:mono} and~\ref{lem:bounds}.
\end{proof}

\begin{lemma}[middle region]\label{lem:middle}
$Q_\varphi(\alpha,x)>0$ for $\alpha\in[4,8.803377]$ and $e^{\psi(1+\alpha)}\le x\le16$.
\end{lemma}
\begin{proof}
Each of the first four ranges is cut into $20$ equal slabs and $[8,8.803377]$ into $16$ ($96$
slabs, of width at most $0.0503$). On a slab with left end
$a$, the $x$-interval runs from a certified lower bound of $e^{\psi(1+a)}$ to $16$. Boxes are
bisected adaptively. On a box the program bounds $T=e^{-x}S_\varphi$ from below by the larger of a
direct ball evaluation and the mean-value form
\[
 T(\alpha_c,x_c)-\sup|\partial_\alpha T|\,r_\alpha-\sup|\partial_xT|\,r_x ,
\]
where $(\alpha_c,x_c)$ is the centre, $r_\alpha$ and $r_x$ are the half-widths, and the suprema are
enclosures obtained by forward-mode differentiation in ball arithmetic. It also bounds
$\alpha^2e^{-x}B_\varphi$ from above over the box. A box is accepted when the resulting lower bound
of $Q_\varphi$ is positive. Otherwise the direction with the larger first-order term is bisected,
and a depth guard of $40$ can only cause failure. The run accepts $413$ boxes in $730$
evaluations, with maximum depth $30$. The smallest certified lower bound is $6.1219\cdot10^{-9}$, near the numerical critical-point candidate.
\end{proof}

\begin{lemma}[witness]\label{lem:witness}
At $\alpha=8.8033777$ and $x=10.7538996$,
\[
 \partial_\alpha U_\alpha(x)
 \in[3.256052269307664\cdot10^{-7}\ \mathbin{\pm}\ 5\cdot10^{-22}]>0.
\]
\end{lemma}
\begin{proof}
Enclose the first $40$ terms of the derivative series in Arb. The absolute
$n$th term for $n\ge41$ is bounded by
\[
 b(n)=n\bigl(\log(1+\alpha n)+|\log x|\bigr)
       \left(\frac{ex}{\alpha n}\right)^{\alpha n}.
\]
This follows from $\Gamma(1+y)\ge(y/e)^y$ and
$0<\psi(1+y)\le\log(1+y)$ in this range. To justify every tail ratio,
set $h(t)=\log(1+\alpha t)+|\log x|$. Since $h>0$, $h'>0$ and $h''<0$,
\[
 (\log b)''=-t^{-2}+h''/h-(h'/h)^2-\alpha/t<0.
\]
Thus $b(n+1)/b(n)$ decreases. The certificate checks
$b(42)/b(41)<1/2$, implying that the entire tail is at most
$2b(41)<7\cdot10^{-392}$. The enclosed partial sum minus this bound is
strictly positive. This calculation is independent of the ray formula.
\end{proof}

\begin{proof}[Proof of Theorem~\ref{thm:sharp}]
(i) On $(0,2)$ use \cite[Theorem~A]{FS}, whose proof gives strict inequalities, and on $[2,4]$
Theorem~\ref{thm:24}. Let
$\alpha\in[4,8.803377]$ and $x>0$. If $x\le e^{\psi(1+\alpha)}$, Lemma~\ref{lem:small} applies. If
$e^{\psi(1+\alpha)}\le x\le16$, Lemmas~\ref{lem:middle} and~\ref{lem:bounds} give
$\alpha^2e^{-x}H_\alpha(x)\ge Q_\varphi(\alpha,x)>0$. If $x\ge16$, use Lemma~\ref{lem:tail16}.
Hence $H_\alpha(x)>0$ on $[2,8.803377]\times(0,\infty)$, and continuity of $\alpha\mapsto U_\alpha(x)$
at $\alpha=2$ gives strict decrease on $(0,8.803377]$.

(ii) This is Lemma~\ref{lem:witness}. By continuity, $\partial_\alpha U_\alpha(10.7538996)>0$ on a
neighbourhood of $8.8033777$, so $U_\alpha(10.7538996)$ increases on an interval to the left of
$8.8033777$. Hence $\alpha^*<8.8033777$; part~(i) gives $\alpha^*\ge8.803377$.
\end{proof}

\subsection{The numerical critical-point candidate}\label{subsec:critical}
Solving $\partial_\alpha U_\alpha(x)=0$ and
$\partial_x\partial_\alpha U_\alpha(x)=0$ numerically gives
\[
 \widehat\alpha=8.80337769817047\ldots,\qquad
 \widehat x=10.7538996236407\ldots.
\]
These are numerical candidates. The computation does not certify uniqueness,
the equality $\widehat\alpha=\alpha^*$, or a nondegenerate first fold.
At an exact zero of $H$ the contour formula requires
$e^x+S_\varphi(\alpha,x)+\alpha^2J_\varphi(\alpha,x)=0$;
the ray remainder must not be omitted. The certified transition interval is
the result of Theorem~\ref{thm:sharp}. The new global theorem in
Section~\ref{sec:no-recovery} replaces finite-grid observations of failure
above that window by a proof on an unbounded interval.

\section{Proofs of the corollaries}\label{sec:cor}
\begin{proof}[Proof of Corollary~\ref{cor:beta}]
For $\beta=1$ this is Theorem~\ref{thm:sharp}(i). For $\beta>1$, the positive series and
$B(1+\alpha n,\beta-1)=\Gamma(1+\alpha n)\Gamma(\beta-1)/\Gamma(\beta+\alpha n)$ give
\[
 E_{\alpha,\beta}(x^\alpha)=\frac{1}{\Gamma(\beta-1)}\int_0^1(1-t)^{\beta-2}E_\alpha\bigl((tx)^\alpha\bigr)\,dt .
\]
If $0<\alpha_1<\alpha_2\le8.803377$, the integrand for $\alpha_1$ is larger than that for $\alpha_2$
at every $t\in(0,1)$ by Theorem~\ref{thm:sharp}(i), and integration preserves the strict
inequality. This argument is taken from the draft, where it was used on $(0,4]$.
\end{proof}

\begin{proof}[Proof of Corollary~\ref{cor:bounds}]
By Theorem~\ref{thm:sharp}(i), $U_2>U_\alpha>U_4$ for $2<\alpha<4$, and $U_4>U_\alpha>U_8$ for
$4<\alpha<8$. For an integer $n\ge1$, summing the exponential series over residue classes modulo
$n$ gives $U_n(x)=\frac1n\sum_{k=0}^{n-1}\exp(xe^{2\pi ik/n})$. For $n=2,4,8$ this yields $\cosh x$,
$\frac12(\cosh x+\cos x)$ and
$\frac14\bigl(\cosh x+\cos x+2\cos\frac{x}{\sqrt2}\cosh\frac{x}{\sqrt2}\bigr)$.
\end{proof}

\section{A global no-recovery theorem}
\label{sec:no-recovery}
The input manuscript sampled orders only up to $30$ when discussing
whether monotonicity might reappear at larger orders. The next theorem
replaces that finite sampling statement by a result on an unbounded
parameter interval. It is an additional theorem, not a claim to resolve
an independently documented classical conjecture.

\begin{theorem}[No recovery above the certified transition window]
\label{thm:no-recovery}
For every real $\alpha\ge 8.8033777$, there exists $x>0$ such that
$\partial_\alpha E_\alpha(x^\alpha)>0$. More explicitly,
\begin{enumerate}
\item for $8.8033777\le\alpha\le9$, one can take $x=10.7538996$;
\item for every $\alpha\ge9$, one can take $x=5\alpha/4$, and
\[
 \left.\partial_\alpha E_\alpha(x^\alpha)\right|_{x=5\alpha/4}
 >\frac{2459}{10^6}\,
 \frac{(5\alpha/4)^\alpha}{\Gamma(1+\alpha)}.
\]
\end{enumerate}
The derivative in the second assertion is a partial derivative at fixed
$x$, evaluated afterwards at $x=5\alpha/4$; it is not the total derivative
of $\alpha\mapsto E_\alpha((5\alpha/4)^\alpha)$.
\end{theorem}

\begin{lemma}[Gamma estimates]
\label{lem:gamma-no-recovery}
For every $z>0$,
\[
 \sqrt{2\pi z}(z/e)^z<\Gamma(1+z)
 <\sqrt{2\pi z}(z/e)^z e^{1/(12z)},
 \qquad \psi(1+z)<\log z+\frac1{2z}.
\]
\end{lemma}
\begin{proof}
Binet's formula \cite[5.9.10]{DLMF} writes
\[
 \log\Gamma(1+z)=(z+\tfrac12)\log z-z+
 \tfrac12\log(2\pi)+R(z),\qquad
 R(z)=2\int_0^\infty\frac{\arctan(t/z)}{e^{2\pi t}-1}\,dt.
\]
Here $0<R(z)<1/(12z)$, by $0<\arctan(t/z)<t/z$ and
$\int_0^\infty t/(e^{2\pi t}-1)\,dt=1/24$.
Differentiation under the integral gives $R'(z)<0$.
Exponentiating and differentiating the identity gives the claims.
\end{proof}

\begin{proof}[Proof of Theorem~\ref{thm:no-recovery} for $\alpha\ge9$]
Put $c=5/4$, $x=c\alpha$, and
\[
 b_n=\frac{(ec)^{n-1}}{n^n},\qquad
 r_n=\frac{x^{\alpha(n-1)}\Gamma(1+\alpha)}{
                    \Gamma(1+n\alpha)}.
\]
Lemma~\ref{lem:gamma-no-recovery} gives, for $n\ge2$,
\[
 r_n<\frac{e^{1/(12\alpha)}}{\sqrt n}\,b_n^\alpha.
\]
Moreover $b_2=5e/16<1$, and
\[
 \frac{b_{n+1}}{b_n}
 =\frac{ec}{n+1}\left(\frac{n}{n+1}\right)^n
 \le \frac{ec}{2(n+1)}<1
 \quad(n\ge2),
\]
since $(1+1/n)^n\ge2$ and $e<3$.
In particular $0<b_n<1$ for every $n\ge2$.
The derivative series, divided by its positive first amplitude, satisfies
\begin{align*}
 \frac{\Gamma(1+\alpha)}{x^\alpha}\partial_\alpha U_\alpha(x)
 &\ge \log c-\frac1{2\alpha}
 -\sum_{n=2}^\infty n r_n
       \left(\log\frac n c+\frac1{2n\alpha}\right)\\
 &\ge\log c-\frac1{18}-e^{1/108}\sum_{n=2}^\infty h_n,
 \qquad
 h_n=\sqrt n\,b_n^9
           \left(\log\frac n c+\frac1{18n}\right).
\end{align*}
All quantities in the last sum are positive. For $n\ge3$,
$c(n+1)\le n^2$ implies
$\log((n+1)/c)\le2\log(n/c)$, and hence
\[
 \frac{h_{n+1}}{h_n}
 <4\left(\frac{15}{32}\right)^9<\frac12.
\]
This estimates the entire infinite tail, giving
$\sum_{n\ge2}h_n<h_2+2h_3$. Therefore the normalized derivative is
bounded below by the explicit constant
\begin{align*}
 \delta={}&\log\frac54-\frac1{18}\\
 &-e^{1/108}\left[
 \sqrt2\left(\frac{5e}{16}\right)^9
        \left(\log\frac85+\frac1{36}\right)
 +2\sqrt3\left(\frac{25e^2}{432}\right)^9
        \left(\log\frac{12}5+\frac1{54}\right)\right].
\end{align*}
A scalar Arb evaluation at 128 bits encloses this constant as
\[
 \delta\in[0.0024592813467114626274414363395958132
                  \mathbin{\pm}6.67\cdot10^{-38}],
\]
and certifies $\delta>2459/10^6$. The accompanying listing makes this
single numerical inequality reproducible. This proves the second part.
\end{proof}

\begin{proof}[Proof for the finite parameter interval]
Fix $x_0=10.7538996$. Let $D_{20}(\alpha)$ be the first 20 terms of
\[
 D(\alpha)=\partial_\alpha U_\alpha(x_0)
 =\sum_{n=1}^\infty
 \frac{n x_0^{\alpha n}}{\Gamma(1+\alpha n)}
       \bigl(\log x_0-\psi(1+\alpha n)\bigr).
\]
A simple uniform tail bound avoids any sampling inference. For
$8\le\alpha\le9$ and $n\ge21$, using $1<x_0<11$, $e<3$,
$0<\psi(1+\alpha n)<\log(\alpha n)+1/(2\alpha n)<9n+1$, and
Lemma~\ref{lem:gamma-no-recovery}, the absolute value of the $n$th
summand is at most
\[
 R_n=13n^2\left(\frac{33}{8n}\right)^{8n}.
\]
Indeed $33/(8n)<1$, so replacing $\alpha n$ by $8n$ in both the
base bound and exponent only increases this upper bound. Furthermore,
\[
 \frac{R_{n+1}}{R_n}
 =\left(\frac{n+1}{n}\right)^2
  \left(\frac{33}{8(n+1)}\right)^8
  \left(\frac n{n+1}\right)^{8n}
 \le4\left(\frac{33}{8(n+1)}\right)^8<\frac12
 \quad(n\ge21).
\]
Consequently
\[
 |D-D_{20}|\le2R_{21}
 =26\cdot21^2\left(\frac{33}{168}\right)^{168}
 <2.079\cdot10^{-115}=:\varepsilon.
\]
The code uses the exact expression $2R_{21}$, not the displayed rounded
upper bound. It starts from the exact rational interval
$[88033777/10^7,9]$ and evaluates $D_{20}-2R_{21}$ on each interval
using 128-bit Arb arithmetic. An interval is accepted only if its entire
enclosure is strictly positive; otherwise it is bisected at its exact
rational midpoint. The depth guard can only stop the proof with failure.
The verified run accepts 1544 intervals in 3087 evaluations, with maximum
depth 33. The smallest accepted lower bound is greater than
$1.5942\cdot10^{-9}$. Every split preserves the union of the current
intervals, and the stack empties, so the certificate covers every real
$\alpha$ in the stated interval. Thus $D(\alpha)>0$ throughout it.
\end{proof}

\begin{remark}
Combined with a lower monotonicity certificate, this theorem leaves only
the stated narrow transition interval undecided. It does not identify an
exact critical order, certify the uniqueness of a critical point, or
exclude isolated behavior inside that transition interval. The
unbounded conclusion is stronger than any finite grid search, but it
should not be relabeled as the solution of an unrelated named conjecture.
\end{remark}

\section{A strict hierarchy of two-parameter thresholds}
\label{sec:beta-hierarchy}
This section gives structural consequences beyond the inherited lower bound for
$\beta\geq1$. These are proved here; no claim is made that they resolve a previously
published conjecture other than the one identified elsewhere in this paper.
Put
\[
 F_{\alpha,\beta}(x)=E_{\alpha,\beta}(x^\alpha),\qquad
 H_{\alpha,\beta}(x)=-\partial_\alpha F_{\alpha,\beta}(x),
\]
and, for $\beta\geq1$, define
\[
 A(\beta)=\sup\{a>0: H_{\alpha,\beta}(x)\geq0
                 \text{ for all }0<\alpha\leq a,\ x>0\}.
\]
The derivative formulation is equivalent to non-increase on the indicated
$\alpha$-interval. In particular $A(1)=\alpha^*$.

\begin{lemma}[Uniform positivity at the two ends]\label{lem:beta-ends}
Fix $\beta\geq1$ and a compact interval $[a,b]\subset[2,\infty)$.
There are $0<r<R<\infty$ such that
$H_{\alpha,\beta}(x)>0$ for every $\alpha\in[a,b]$ and
$x\in(0,r]\cup[R,\infty)$. More precisely, for some $\rho\in(0,1)$,
\[
 x^{\beta-1}H_{\alpha,\beta}(x)
       =\frac{e^x}{\alpha^2}+O(e^{\rho x})
       \qquad (x\to\infty),
\]
with the error uniform in $\alpha\in[a,b]$.
\end{lemma}
\begin{proof}
Termwise differentiation gives
\[
 H_{\alpha,\beta}(x)=\sum_{n\geq1}
 \frac{n x^{\alpha n}}{\Gamma(\beta+\alpha n)}
       \bigl(\psi(\beta+\alpha n)-\log x\bigr).
\]
Thus $r=\exp(\psi(\beta+a))$ suffices at the small end; at $x=r$
the terms with $n\geq2$ are strictly positive.
For the other end the Laplace transform is
\[
 \int_0^\infty e^{-sx}x^{\beta-1}H_{\alpha,\beta}(x)\,dx
   =\frac{s^{\alpha-\beta}\Log s}{(s^\alpha-1)^2},
       \qquad \Re s>1.
\]
Interchange and inversion follow by the same gamma-moment and absolute-integrability
argument used for $\beta=1$; on a vertical line the transform decays like
$O(|s|^{-\alpha-\beta}\log|s|)$ uniformly on the compact parameter interval.
Choose
\[
 \max\{0,\cos(2\pi/b)\}<\rho<1.
\]
Every nonprincipal root of $s^\alpha=1$ has real part at most
$\cos(2\pi/b)$, so shifting the inversion line to $\Re s=\rho$
crosses only $s=1$. Its residue after multiplication by $e^{xs}$ is
$e^x/\alpha^2$. On the new line the denominator is bounded away from zero
uniformly on bounded imaginary parts and $\alpha\in[a,b]$; the displayed decay
controls the remaining tails uniformly. The remainder is therefore bounded
in absolute value by $M e^{\rho x}$ for a finite constant $M$ independent of
$\alpha$. Taking $e^{(1-\rho)R}>b^2M$ proves the claim.
\end{proof}

\begin{lemma}[Finiteness of every threshold]\label{lem:beta-finite}
For every fixed $\beta\geq1$ and every fixed $c\in(1,4/e)$,
\[
 \partial_\alpha F_{\alpha,\beta}(c\alpha)>0
 \quad\text{for all sufficiently large }\alpha.
\]
The derivative is taken at fixed $x$, then evaluated at $x=c\alpha$;
it is not the derivative of the composed function
$\alpha\mapsto F_{\alpha,\beta}(c\alpha)$.
Consequently $A(\beta)<\infty$.
\end{lemma}
\begin{proof}
For a fixed large $\alpha$ let
$t_n=(c\alpha)^{\alpha n}/\Gamma(\beta+\alpha n)$.
Then
\[
 \partial_\alpha F_{\alpha,\beta}(c\alpha)
 =\sum_{n\geq1}n t_n
        \bigl(\log(c\alpha)-\psi(\beta+\alpha n)\bigr).
\]
For $\alpha\geq1$ and $y>0$, Jensen's inequality applied to a gamma random
variable $G_y$ of unit rate gives
$\Gamma(y+\alpha)/\Gamma(y)=\mathbb E G_y^\alpha\geq y^\alpha$.
It follows that, for $n\geq2$,
\[
 \frac{t_{n+1}}{t_n}\leq
 \left(\frac{c\alpha}{\beta+n\alpha}\right)^\alpha
 \leq(c/n)^\alpha\leq(c/2)^\alpha.
\]
For all sufficiently large $\alpha$ this last bound is at most $1/2$.
The elementary bounds
$\log(y-1)\leq\psi(y)\leq\log y$ for $y>1$ follow respectively
from Jensen applied to $G_y^{-1}$ and $G_y$.
Hence, for $\alpha\geq1$, $n\geq1$, and $\beta\geq1$,
\[
 \bigl|\log(c\alpha)-\psi(\beta+n\alpha)\bigr|
 \leq\log c+\log(n+\beta).
\]
Consequently the absolute sum of the terms with $n\geq2$ is at most
\[
 C_{\beta,c}t_2,\qquad
 C_{\beta,c}=\sum_{n\geq2}n\bigl(\log c+\log(n+\beta)\bigr)2^{-(n-2)}<\infty.
\]
Stirling's formula with $\beta$ fixed gives
\[
 \frac{t_2}{t_1}
  =2^{1/2-\beta}\left(\frac{ce}{4}\right)^\alpha(1+O(\alpha^{-1}))
  \longrightarrow0,
\]
while $\log(c\alpha)-\psi(\beta+\alpha)\to\log c>0$.
The positive $n=1$ contribution therefore dominates the whole tail.
\end{proof}

\begin{theorem}[Strict smoothing of the monotonicity boundary]
\label{thm:beta-hierarchy}
For every $\beta\geq1$, $A(\beta)$ is finite and the defining inequalities
hold at $\alpha=A(\beta)$. Moreover:
\begin{enumerate}
\item If $1\leq\beta_1<\beta_2$, then $A(\beta_1)<A(\beta_2)$.
In particular $A(\beta)>\alpha^*$ for every $\beta>1$.
\item The function $\beta\mapsto A(\beta)$ is right-continuous on $[1,\infty)$.
\item For every $\beta\geq1$ there exists $x_\beta\in(0,\infty)$ such that
\[
 H_{A(\beta),\beta}(x_\beta)=0,
 \qquad
 \partial_xH_{A(\beta),\beta}(x_\beta)=0.
\]
\end{enumerate}
Thus the first loss of parameter monotonicity occurs at a finite interior
contact point, and increasing the second parameter strictly postpones it.
This theorem does not assert uniqueness or nondegeneracy of the contact.
\end{theorem}
\begin{proof}
The already established one-parameter result and positive beta integration
show that $A(\beta)\geq8.803377>2$. Finiteness follows from
Lemma~\ref{lem:beta-finite}. The set of admissible endpoints is an initial
interval. Continuity of $H$ implies that its finite supremum is admissible.

For $\delta=\beta_2-\beta_1>0$ the beta integral gives the exact identity
\[
 H_{\alpha,\beta_2}(x)=\frac1{\Gamma(\delta)}
 \int_0^1 t^{\beta_1-1}(1-t)^{\delta-1}
                    H_{\alpha,\beta_1}(tx)\,dt.
\]
This follows term by term; all relevant derivatives converge locally
uniformly, and the beta weights are integrable at both endpoints. Near $t=0$,
$H_{\alpha,\beta_1}(tx)=O(t^{a_0}|\log t|)$ uniformly on any compact
$\alpha$-neighborhood with positive lower endpoint $a_0$, providing an
integrable majorant for the differentiation.
At $\alpha=A(\beta_1)$ the integrand is nonnegative and it is strictly
positive for all sufficiently small $t>0$. Therefore
$H_{A(\beta_1),\beta_2}(x)>0$ for every $x>0$.
Apply Lemma~\ref{lem:beta-ends} to a compact neighborhood of $A(\beta_1)$.
It supplies uniform positivity outside a compact $x$-interval. On that
compact interval strict positivity at $A(\beta_1)$ and continuity supply
positivity for $\alpha\in[A(\beta_1),A(\beta_1)+\varepsilon]$ with some
$\varepsilon>0$. At every smaller positive $\alpha$ nonnegativity follows
from the same beta integral. Thus
$A(\beta_2)\geq A(\beta_1)+\varepsilon$.

For right-continuity fix $\beta$ and any $a>A(\beta)$. By nonadmissibility
there exist $0<\alpha\leq a$ and $x>0$ with
$H_{\alpha,\beta}(x)<0$. Continuity in $\beta$ preserves this strict
inequality in a neighborhood of $\beta$, giving
$A(\beta')<\alpha\leq a$ there. Hence
$\limsup_{\beta'\to\beta}A(\beta')\leq A(\beta)$.
Monotonicity gives the reverse inequality when $\beta'\downarrow\beta$.

Finally set $A=A(\beta)$. For each positive integer $j$ there exist
$\alpha_j\in(A,A+1/j]$ and $x_j>0$ with
$H_{\alpha_j,\beta}(x_j)<0$.
Lemma~\ref{lem:beta-ends} confines all sufficiently late $x_j$ to a single
compact interval $[r,R]\subset(0,\infty)$. A convergent subsequence has
limit $x_\beta$. Continuity gives $H_{A,\beta}(x_\beta)\leq0$, whereas
admissibility gives the opposite inequality. Since $H_{A,\beta}$ is
nonnegative and differentiable on $(0,\infty)$, its derivative vanishes at
this interior zero.
\end{proof}

\begin{remark}
The conclusion $A(\beta)>\alpha^*$ is qualitative. It does not supply a
numerical lower bound larger than $8.8033777$ for any specific $\beta$ without
additional estimates or certified computation. The proof also does not
show that $A(\beta)\to\infty$ as $\beta\to\infty$, nor does it establish
left-continuity, differentiability, or convexity of $A$.
\end{remark}

\section{A certified extension to order nine for the second parameter}
\label{sec:beta2-certificate}
The preceding hierarchy is qualitative. The following separate computation
shows that its extension is already visible at the concrete value $\beta=2$.
\begin{theorem}\label{thm:beta2-nine}
For every $\beta\geq2$ and $x>0$, the function
$\alpha\mapsto E_{\alpha,\beta}(x^\alpha)$ is strictly decreasing on $(0,9]$.
Moreover $A(\beta)>9$ for every $\beta\geq2$.
\end{theorem}
\begin{proof}
It suffices first to prove $H_{\alpha,2}(x)>0$ for
$4\leq\alpha\leq9$ and $x>0$. The smaller orders already follow from
the one-parameter theorem and the beta integral.
The Laplace transform of $xH_{\alpha,2}(x)$ is $s^{-1}K_\alpha(s)$.
Moving the contour to the same rays as before yields
\[
 xH_{\alpha,2}(x)
 =\frac{e^x+S_{2,\varphi}(\alpha,x)}{\alpha^2}
   +J_{2,\varphi}(\alpha,x),
\]
where, with $p_k=x\sin u_k$ and $u_k=2\pi k/\alpha$,
\[
 S_{2,\varphi}
 =2\sum_{k=1}^{K_\varphi} e^{x\cos u_k}
   \{\cos(p_k-u_k)+u_k\sin(p_k-u_k)-u_kx\sin p_k\}.
\]
Indeed the residue at a nonprincipal pole $w=e^{iu}$ is
\[
 \frac{e^{xw}}{\alpha^2}
       \{w^{-1}(1-iu)+iux\}.
\]
The ray integral is the earlier integral with $r^{\alpha-1}$ replaced
by $r^{\alpha-2}$ and $e^{i\alpha\varphi}$ replaced by
$e^{i(\alpha-1)\varphi}$. Its absolute value is bounded by
\[
 B_{2,\varphi}=
 \frac1{\pi m^2}\left[
 \frac{\Gamma(\alpha-1)}{\lambda^{\alpha-1}}
 \{\varphi+(\psi(\alpha-1)-\log\lambda)_+\}
 +2\min\left\{\frac1{(\alpha-1)^2},
       \frac{\Gamma(\alpha-2)}{e\lambda^{\alpha-2}}\right\}\right],
\]
where $\lambda=-x\cos\varphi$ and $m^2$ retains its earlier definition.
This follows by the same integral estimates, now with exponent
$r^{\alpha-2}$; here $\alpha\geq4$, so every gamma integral converges.
Consequently
\[
 \alpha^2e^{-x}xH_{\alpha,2}(x)
 \geq Q_2:=1+e^{-x}S_{2,\varphi}-\alpha^2e^{-x}B_{2,\varphi}.
\]
A simpler sufficient condition is $G_2<1$, where
\[
 G_2=\sum_{k=1}^{K_\varphi}
       2e^{-x(1-\cos u_k)}(1+u_k+u_kx)
          +\alpha^2e^{-x}B_{2,\varphi}.
\]
For $\alpha\leq9$, $u_k\geq2\pi/9$ and $u_k<\pi$.
Since $(1-\cos(2\pi/9))^{-1}<5$, each residue bound decreases
strictly for $x\geq16$. The ray bound also decreases there, term by
term, exactly as in the one-parameter argument. Thus $G_2$ decreases
for $x\geq16$.

Use the five parameter ranges and ray angles
\[
\begin{array}{c|ccccc}
\alpha&[4,5]&[5,6]&[6,7]&[7,8]&[8,9]\\
\varphi/\pi&13/20&9/10&19/25&33/50&83/100\\
K_\varphi&1&2&2&2&3.
\end{array}
\]
The attached reproducible Arb listing checks the pole geometry and
$m^2>1/2$ on each range. It then establishes the following facts in
128-bit ball arithmetic (python-flint 0.9.0):
\begin{enumerate}
\item On 1,000 rational parameter intervals,
$G_2(\alpha,16)<0.610$; the largest computed upper bound is
$0.609741652806147591445611468801\ldots$.
\item On the middle region
$e^{\psi(\alpha+2)}\leq x\leq16$, adaptive bisection accepts
231 boxes in 362 evaluations, with maximum depth 3 and
$Q_2>0.00134654$ on every accepted box.
\end{enumerate}
For the middle region, each of 100 exact rational parameter slabs
starts at a certified lower bound of $e^{\psi(a+2)}$, where $a$ is
its left endpoint. The lower enclosure of $e^{-x}S_{2,\varphi}$
uses the larger of a direct interval bound and its centered
mean-value bound, with both partial derivatives enclosed by automatic
differentiation. Each rejected box is replaced by its two halves;
the empty final stack therefore certifies full coverage.
For $0<x\leq e^{\psi(\alpha+2)}$, every derivative-series term
is nonnegative and those with $n\geq2$ are positive.
These three regions prove $H_{\alpha,2}(x)>0$ on $[4,9]\times(0,\infty)$.

The beta integral with $\beta_1=2$ transfers the resulting strict
order inequalities to every $\beta>2$. This proves strict decrease
on $(0,9]$ for every $\beta\geq2$. In addition, the same integral
gives $H_{9,\beta}(x)>0$ for every $x>0$. Uniform positivity at the
two ends and compactness of the remaining $x$-interval then extend
nonnegativity to a common interval $[9,9+\varepsilon]$, for some
$\varepsilon>0$ depending on $\beta$. Thus $A(\beta)>9$.
\end{proof}

\begin{remark}
This explicit interval extends beyond the certified obstruction
$\alpha^*<8.8033777$ at $\beta=1$. It proves an actual distinction
between the parameter ranges; it does not determine $A(2)$ exactly.
\end{remark}

\section{Transfer to complete monotonicity and fractional equations}
\label{sec:transfers}
The preceding theorems also answer two explicit extension questions: which
elementary Laplace kernels inherit infinitely many derivative inequalities,
and which fractional integral equations inherit order comparison?

\begin{theorem}[An exact complete-monotonicity criterion]\label{thm:CM}
For $\alpha\ge2$ and $\beta\ge1$, put
\[
 L_{\alpha,\beta}(s)=\frac{s^{\alpha-\beta}\log s}{(s^\alpha-1)^2},\qquad s>1.
\]
Then $L_{\alpha,\beta}$ is completely monotone on $(1,\infty)$ if and only if
$H_{\alpha,\beta}(x)\ge0$ for every $x>0$. Here complete monotonicity means
$(-1)^m L^{(m)}(s)\ge0$ for every integer $m\ge0$ and $s>1$.
In particular these inequalities hold for every $2\le\alpha\le8.803377$
and $\beta\ge1$, whereas $L_{\alpha,1}$ is not completely monotone for any
$\alpha\ge8.8033777$.
\end{theorem}
\begin{proof}
Set $h(x)=x^{\beta-1}H_{\alpha,\beta}(x)$. Its Laplace transform is $L$, by
Lemma~\ref{lem:beta-ends}. The series at zero shows that $h$ extends continuously
to zero with value zero; the contour bound at infinity gives $|h(x)|\le Ce^{cx}$
for every fixed $c>1$, with a suitable $C$. Differentiation under the integral gives
\[
 (-1)^m L^{(m)}(s)=\int_0^\infty x^m e^{-sx}h(x)\,dx.
\]
Thus $h\ge0$ implies complete monotonicity. Conversely fix $t>0$ and let
$Y_m$ have gamma shape $m+1$ and rate $m/t$. For $m/t>1$,
\[
 \mathbb E h(Y_m)=\frac{(m/t)^{m+1}}{m!}
                  (-1)^mL^{(m)}(m/t)\ge0.
\]
Its mean tends to $t$ and variance tends to zero. Moreover for every fixed
$c>1$, $\mathbb E e^{2cY_m}=(1-2ct/m)^{-(m+1)}$ is bounded for all sufficiently
large $m$. The exponential bound on $h$, Cauchy--Schwarz on the complement of
a neighborhood of $t$, and continuity inside that neighborhood imply
$\mathbb E h(Y_m)\to h(t)$. Therefore $h(t)\ge0$. The parameter ranges follow
from Theorems~\ref{thm:sharp} and~\ref{thm:no-recovery} and beta integration.
\end{proof}

This equivalence transfers the same transition to an elementary function of
$s$ with infinitely many derivative constraints. It is not a statement about
the different reciprocal $1/E_{\alpha,\beta}$, nor about the positions of
complex zeros. In particular the complex-modulus conjectures formulated in
\cite[Section 4]{GGPS} do not follow from Theorem~\ref{thm:CM}.

\begin{theorem}[Comparison for nondecreasing forcing]\label{thm:volterra}
Let $g:[0,\infty)\to[0,\infty)$ be continuously differentiable with $g'\ge0$,
and fix $\lambda>0$. For $\alpha>0$ let $u_\alpha$ be the unique continuous
solution of
\[
 u_\alpha(t)=g(t)+\lambda^\alpha I^\alpha u_\alpha(t),\qquad
 I^\alpha f(t)=\frac1{\Gamma(\alpha)}\int_0^t(t-s)^{\alpha-1}f(s)\,ds.
\]
If $0<\alpha_1<\alpha_2\le8.803377$, then
$u_{\alpha_1}(t)\ge u_{\alpha_2}(t)$ for every $t\ge0$.
For every $t>0$ the inequality is strict if $g(0)>0$ or
$\int_0^t g'(s)\,ds>0$.
The maximal initial interval of orders on which this non-strict comparison
holds for every such $g$, every $\lambda>0$ and every time is exactly
$(0,\alpha^*]$.
\end{theorem}
\begin{proof}
The locally convergent resolvent series
$u_\alpha=\sum_{n\ge0}\lambda^{\alpha n}I^{\alpha n}g$ gives existence.
Uniqueness follows by iterating the homogeneous equation on $[0,T]$ and using
$\|I^{n\alpha}f\|_\infty\le T^{n\alpha}\|f\|_\infty/\Gamma(1+n\alpha)$.
Termwise convolution or integration by parts gives
\[
 u_\alpha(t)=g(0)U_\alpha(\lambda t)
          +\int_0^tU_\alpha(\lambda(t-s))g'(s)\,ds.
\]
Every weight is nonnegative, so Theorem~\ref{thm:sharp} proves comparison and
strictness. The same argument works on the whole defining interval for
$\alpha^*$, including its endpoint by continuity. Conversely the admissible
test input $g\equiv1$ gives $u_\alpha(t)=U_\alpha(\lambda t)$; any universal
comparison interval must therefore lie inside the defining one for $\alpha^*$.
\end{proof}

\section{Assessment and remaining problems}\label{sec:assessment}
The input manuscript contains a substantial computer-assisted argument in
special-function inequalities. Its main claim addresses the explicit
expectation in Ferreira--Simon, Remark~1(a). This revision independently
reproduces its three supplied certificates, repairs the witness-tail argument,
corrects a rounded certificate bound, and separates numerical observations
from proved statements. The core proof combines global complex analysis with
finite interval certificates; it is not a finite sample extrapolation.

The additional results have different roles. The no-recovery theorem converts
a finite numerical observation into an unbounded theorem. The strict hierarchy
$A(\beta_2)>A(\beta_1)$ describes how beta integration changes the actual
boundary, rather than merely preserving an already certified interval. The
complete-monotonicity equivalence and Volterra theorem are rigorously proved
transfers of that boundary. These are theorems proved in this revision, but
historical novelty and priority are not established by the limited literature
search. We do not claim the resolution of an additional independently named
long-standing conjecture.

What remains is specific: identify the exact first contact inside the certified
window; prove its uniqueness and nondegeneracy, if true; determine whether
$A(\beta)$ is left-continuous or differentiable; and determine its growth as
$\beta\to\infty$. The value $8.80337769817047\ldots$ remains a numerical candidate.
No theorem here proves irrationality, transcendence, or a statement about
unrelated classical constants. Independent mathematical refereeing and a
formal proof-assistant check have not taken place.

\section{Verification record}\label{sec:record}
All original supplied programs were extracted and rerun successfully with
python-flint 0.9.0. This revision uses the same analytic estimates and strengthens
the last range to $8803377/10^6$. The witness endpoint is $88033777/10^7$.
The standalone listings below are part of the single TeX source.
\begin{center}\small
\begin{tabular}{p{4.2cm}p{9.5cm}}\toprule
Certificate & Reproduced or strengthened result\\\midrule
Proof A, 96 bits & 16,749 rectangles; 33,497 evaluations; depth 16; exact area 27.\\
Proof B, 128 bits & 2,000 intervals; upper bounds $0.199825$ and $0.182407$.\\
Transition window, 128 bits & 413 accepted boxes; 730 evaluations; depth 30;
minimum $Q>6.1219\cdot10^{-9}$; tail $G<0.501$.\\
Positive derivative witness & $D(8.8033777,10.7538996)>3.2560\cdot10^{-7}$;
tail $<7\cdot10^{-392}$.\\
No recovery, 128 bits & 1,544 intervals; 3,087 evaluations; depth 33;
minimum $D>1.5942\cdot10^{-9}$.\\
$\beta=2$, 128 bits & 231 boxes; 362 evaluations; depth 3;
minimum $Q>0.0013465$; tail $G<0.610$.\\\bottomrule
\end{tabular}
\end{center}
Each interval program starts from exactly specified endpoints and accepts a box
only after an interval inequality succeeds. Bisection preserves coverage;
a depth guard causes failure, never acceptance. Infinite ranges and infinite
series tails are controlled analytically. Floating-point samples are not used
to infer a universal sign. The new scalar margin for $\alpha\ge9$ is
$\delta>0.002459$, proved by one Arb enclosure and the analytic tail estimates.
Optional \texttt{--numeric} routines in the inherited listings are diagnostics
only and do not enter the proof.

To reproduce, install \texttt{python-flint==0.9.0}, save the five Python
listings under their stated names, and execute them with Python 3.
The optional diagnostics additionally require \texttt{mpmath}.

\begin{thebibliography}{9}
\bibitem{FS}
R.~A.~C. Ferreira and T.~Simon, \emph{Mittag--Leffler functions and convex ordering},
arXiv:2512.04940v1, 4 December 2025.
\href{https://arxiv.org/abs/2512.04940v1}{arxiv.org/abs/2512.04940v1}.
Theorem~A and Remark~1(a).

\bibitem{GKMR}
R.~Gorenflo, A.~A. Kilbas, F.~Mainardi and S.~V. Rogosin,
\emph{Mittag--Leffler Functions, Related Topics and Applications}, 2nd ed.,
Springer Monographs in Mathematics, Springer, Berlin, 2020.

\bibitem{Arb}
F.~Johansson, \emph{Arb: efficient arbitrary-precision midpoint-radius interval arithmetic},
IEEE Transactions on Computers 66 (2017), 1281--1292.
\href{https://doi.org/10.1109/TC.2017.2690633}{doi:10.1109/TC.2017.2690633}.

\bibitem{flint}
The python-flint developers, \emph{python-flint 0.9.0}, Python interface to FLINT and Arb.
\href{https://python-flint.readthedocs.io/}{python-flint.readthedocs.io}.

\bibitem{DLMF} NIST Digital Library of Mathematical Functions,
Section 5.9, Binet's formula (5.9.10).
\href{https://dlmf.nist.gov/5.9.E10}{dlmf.nist.gov/5.9.E10}.
\bibitem{GGPS} R. Garrappa, S. Gerhold, M. Popolizio and T. Simon,
\emph{On some inequalities for the two-parameter Mittag--Leffler function
in the complex plane}, J. Math. Anal. Appl. 551 (2025), 129588.
\href{https://arxiv.org/abs/2410.11852v4}{arXiv:2410.11852v4}.
\end{thebibliography}

\appendix
\section{Programs}\label{app:code}
\subsection{Proof A: compact certificate (\texttt{verify\_bromwich\_certificate.py})}\label{app:A}
\begin{lstlisting}
#!/usr/bin/env python3
"""
Independent re-implementation of the compact certificate of Proof A (Bromwich-line proof),
Lemmas 3.2-3.3 of the paper:

    for every (alpha, x) in [2,4] x [5/2,16]:   e^{-x} H_{alpha,24}(x) > 1/100000,
    H_{alpha,24}(x) = sum_{n=1}^{24} n x^{alpha n} (psi(1 + alpha n) - log x) / Gamma(1 + alpha n),

together with the two scalar inequalities psi(3) > log(5/2) and psi(51) > log(16).

Arb ball arithmetic (python-flint), exact rational box endpoints, adaptive bisection
(alpha is bisected when 8 * width(alpha) > width(x), otherwise x), depth guard = failure.
Coverage invariant: every popped box is either accepted or replaced by two halves.
"""
import sys
import time
from flint import arb, fmpq, ctx

ctx.prec = 96
THRESHOLD = arb(fmpq(1, 100000))


def ball(lo, hi):
    return arb(lo).union(arb(hi))


def lower_bound(alo, ahi, xlo, xhi):
    A = ball(alo, ahi)
    X = ball(xlo, xhi)
    lx = X.log()
    s = arb(0)
    for n in range(1, 25):
        nA = n * A
        s += n * (nA * lx - (1 + nA).lgamma() - X).exp() * ((1 + nA).digamma() - lx)
    return s.lower()


def scalar_inequalities():
    a = arb(3).digamma() - arb(fmpq(5, 2)).log()
    b = arb(51).digamma() - arb(16).log()
    return a, b, (a > 0) and (b > 0)


def certify():
    stack = [(fmpq(2), fmpq(4), fmpq(5, 2), fmpq(16), 0)]
    accepted = evals = maxdepth = 0
    area = fmpq(0)
    minlb = None
    while stack:
        a0, a1, x0, x1, d = stack.pop()
        if d > 60:
            return None
        evals += 1
        lb = lower_bound(a0, a1, x0, x1)
        if lb > THRESHOLD:
            accepted += 1
            maxdepth = max(maxdepth, d)
            area += (a1 - a0) * (x1 - x0)
            if minlb is None or lb < minlb:
                minlb = lb
            continue
        if 8 * (a1 - a0) > (x1 - x0):
            am = (a0 + a1) / 2
            stack += [(a0, am, x0, x1, d + 1), (am, a1, x0, x1, d + 1)]
        else:
            xm = (x0 + x1) / 2
            stack += [(a0, a1, x0, xm, d + 1), (a0, a1, xm, x1, d + 1)]
    return accepted, evals, maxdepth, area, minlb


if __name__ == '__main__':
    t0 = time.time()
    a, b, ok = scalar_inequalities()
    print(f'psi(3) - log(5/2) = {a.str(6)},  psi(51) - log 16 = {b.str(6)}  -> {"OK" if ok else "FAIL"}')
    res = certify()
    if res is None or not ok:
        print('CERTIFICATE FAILED')
        sys.exit(1)
    acc, ev, md, area, mlb = res
    print(f'accepted rectangles = {acc}, evaluations = {ev}, max depth = {md}, '
          f'exact covered area = {area} (target 27), full binary tree: {ev == 2 * acc - 1}')
    print(f'smallest certified lower bound of e^-x H_24 = {mlb.str(5, radius=False)}  (> 1/100000)')
    print(f'elapsed {time.time() - t0:.1f} s')
\end{lstlisting}
\subsection{Proof B: one-dimensional certificate (\texttt{verify\_rays\_2\_4.py})}\label{app:B}
\begin{lstlisting}
#!/usr/bin/env python3
"""
Proof B of Theorem 1.1 (rotated rays), computer-assisted step = Lemma 4.1 of the paper:
for every x > 0, alpha -> E_alpha(x^alpha) is strictly decreasing on [2, 4]
(Ferreira-Simon, arXiv:2512.04940, Remark 1(a)).

Usage
    pip install python-flint mpmath
    python3 verify_rays_2_4.py            # rigorous check (Lemma 4.1), ~0.1 s
    python3 verify_rays_2_4.py --numeric  # also: ray formula (Lemma 2.5) vs series, ~5 s

Notation (all logs/powers principal):  theta = 2*pi/alpha,  lambda = -x*cos(phi),  gamma = alpha*phi,
    m2 = 1 if cos(gamma) <= 0 else sin(gamma)^2,
    B  = [ Gamma(a) lambda^-a (phi + (psi(a) - ln lambda)_+) + 2 min(1/a^2, Gamma(a-1)/(e lambda^(a-1))) ] / (pi m2),
    G  = 1[theta < phi] * 2 exp(-x(1 - cos theta)) (1 + theta x) + a^2 e^-x B.
Lemma 2.6: D(a,x) := d/da E_a(x^a) <= (e^x / a^2) (G - 1).   Lemma 2.7: G is decreasing in x for x >= 1.
Lemma 4.1 (this script): on each sub-interval [a, a'] of [2, 2.5] (phi = 0.7 pi) and [2.5, 4] (phi = 0.9 pi),
    sup_{alpha in [a,a']} G(alpha, X) < 1,  where X is a rigorous lower bound of exp(psi(1 + a)).
All quantities are enclosed with Arb ball arithmetic (python-flint), so rounding errors are accounted for.
"""
import sys
from flint import arb, ctx

ctx.prec = 128
PI = arb.pi()
E = arb(1).exp()


def G_upper(a_lo, a_hi, phi_frac, residue_inside):
    """Rigorous upper bound of sup_{alpha in [a_lo, a_hi]} G(alpha, X), X <= exp(psi(1 + a_lo))."""
    a_lo, a_hi = arb(a_lo), arb(a_hi)
    alpha = arb(((a_lo + a_hi) / 2).mid(), ((a_hi - a_lo) / 2).upper())  # ball containing [a_lo, a_hi]
    assert alpha.contains(a_lo) and alpha.contains(a_hi)
    phi = PI * arb(phi_frac)
    X = (1 + a_lo).digamma().exp().lower()          # X <= e^{psi(1+alpha)} for every alpha in the interval
    assert X > 1                                    # needed for Lemma 2.7
    theta = 2 * PI / alpha

    term_res = arb(0)
    if residue_inside:                              # poles e^{+-i theta} lie inside |arg w| < phi
        assert theta < phi
        term_res = 2 * (-(X * (1 - theta.cos()))).exp() * (1 + theta * X)
    else:                                           # poles lie outside: no residue term
        assert theta > phi

    gam = alpha * phi
    if gam.cos().upper() <= 0:
        m2 = arb(1)
    else:
        m2 = (gam.sin() ** 2).lower()
        assert m2 > 0
    lam = -X * phi.cos()
    assert lam > 0
    lnlam = lam.log()
    p1 = (alpha.gamma().upper() * (-(alpha * lnlam)).exp().upper()
          * (phi.upper() + max((alpha.digamma() - lnlam).upper(), arb(0))))
    lb = min((1 / alpha ** 2).upper(), ((alpha - 1).gamma() / (E * lam ** (alpha - 1))).upper())
    B = (p1 + 2 * lb) / (PI.lower() * m2)
    term_ray = (alpha ** 2).upper() * (-X).exp().upper() * B
    return (term_res.upper() + term_ray.upper()).upper()


def rigorous_check():
    ok = True
    for lo, hi, n, phi_frac, inside, name in [(2, '2.5', 500, '0.7', False, 'Range I  [2,2.5], phi=0.7pi'),
                                               ('2.5', 4, 1500, '0.9', True, 'Range II [2.5,4], phi=0.9pi')]:
        worst, where = arb(0), None
        lo_, hi_ = arb(lo), arb(hi)
        for i in range(n):
            a0 = (lo_ + (hi_ - lo_) * i / n).mid()
            a1 = (lo_ + (hi_ - lo_) * (i + 1) / n).mid()
            g = G_upper(a0, a1, phi_frac, inside)
            if not g < 1:
                ok = False
                print('FAIL', a0, a1, g)
            if g > worst:
                worst, where = g, (a0, a1)
        print(f'{name}: {n} sub-intervals, max upper bound of G = {worst.str(6, radius=False)} '
              f'on [{where[0].str(6, radius=False)}, {where[1].str(6, radius=False)}]')
    print('Lemma 4.1 VERIFIED: G < 1 on all of [2,4]' if ok else 'Lemma 4.1 FAILED')
    return ok


def numeric_contour_check():
    """Non-rigorous sanity check of Lemma 2.5: ray formula vs. term-wise differentiated series."""
    import mpmath as mp
    mp.mp.dps = 40

    def D_series(a, x):
        a, x = mp.mpf(a), mp.mpf(x)
        lx, s, n = mp.log(x), mp.mpf(0), 1
        while True:
            an = a * n
            t = mp.exp(an * lx - mp.loggamma(1 + an))
            s += n * (lx - mp.digamma(1 + an)) * t
            if an > x + 60 and t < mp.mpf(10) ** -50:
                return s
            n += 1

    def D_contour(a, x, phi):
        a, x, phi = mp.mpf(a), mp.mpf(x), mp.mpf(phi)
        th = 2 * mp.pi / a
        res = mp.exp(x) / a ** 2
        if th < phi:
            w = mp.exp(1j * th)
            res += 2 * mp.re(mp.exp(x * w) / a ** 2 * (1 + x * w * 1j * th))
        e, eap = mp.exp(1j * phi), mp.exp(1j * a * phi)
        f = lambda r: mp.exp(x * r * e) * r ** (a - 1) * eap * (mp.log(r) + 1j * phi) / (r ** a * eap - 1) ** 2
        J = mp.im(mp.quad(f, [0, 0.5, 1, 2, 5, mp.inf])) / mp.pi
        return -res - J

    worst = mp.mpf(0)
    for a in ['2.1', '2.3', '2.7', '3.0', '3.4', '3.9']:
        for x in ['0.7', '2.5', '6']:
            ds = D_series(a, x)
            for phi in [0.75 * mp.pi, 0.92 * mp.pi]:
                worst = max(worst, abs(ds - D_contour(a, x, phi)) / abs(ds))
    print(f'Lemma 2.5 numeric check: 36 (alpha, x, phi) cases, max relative difference = {mp.nstr(worst, 3)}')


if __name__ == '__main__':
    ok = rigorous_check()
    if '--numeric' in sys.argv:
        numeric_contour_check()
    sys.exit(0 if ok else 1)
\end{lstlisting}
\subsection{Transition-window certificate (\texttt{verify\_sharp\_range.py})}\label{app:C}
\begin{lstlisting}
#!/usr/bin/env python3
"""
Theorem 1.2 (sharp range), computer-assisted steps = Lemmas 5.1-5.4 of the paper.
Rigorous (Arb ball arithmetic via python-flint) certificate that
    D(alpha, x) := d/dalpha E_alpha(x^alpha) < 0   for all alpha in [4, 8803377/1000000], x > 0,
plus a certified witness D(8.8033777, 10.7538996) > 0.  Hence 8.803377 <= alpha* < 8.8033777.
Usage:  python3 verify_sharp_range.py [--numeric]   (~0.1 s; --numeric adds a ~2 min series cross-check)

Representation used for x >= e^{psi(1+alpha)} (Lemma 2.5; cross-checked against the series by --numeric):
    D = -(e^x + S(alpha, x)) / alpha^2 - J_phi(alpha, x),
    S = sum_{k=1}^{K} 2 e^{x cos u_k} [cos(x sin u_k) - u_k x sin(x sin u_k + u_k)],  u_k = 2 pi k / alpha,
    |J_phi| <= B(alpha, x)  (Lemma 2.6; m^2 > 1/2 asserted on every range below = Lemma 5.1).
Hence D < 0 whenever Q := 1 + e^{-x} S - alpha^2 e^{-x} B > 0.

Regions:  x <= e^{psi(1+alpha)}: termwise sign (no computation).
          e^{psi(1+alpha)} <= x <= 16: 2-D adaptive certificate of Q > 0 (Lemma 5.3).
          x >= 16: analytic bound G = sum_k 2 e^{-x(1-cos u_k)}(1 + u_k x) + alpha^2 e^{-x} B < 1,
                   decreasing in x for x >= 16 (Lemma 2.7); certified at x = 16 on sub-intervals (Lemma 5.2).
Witness:  series partial sum + explicit tail bound, independent of Lemma 2.5 (Lemma 5.4).
"""
import sys
import time
from flint import arb, fmpq, ctx

ctx.prec = 128
PI = arb.pi()
E = arb(1).exp()
X_TOP = arb(16)

# (alpha_lo, alpha_hi, phi / pi, K = number of residue pairs inside |arg w| < phi)
RANGES = [
    (fmpq(4), fmpq(5), fmpq(13, 20), 1),
    (fmpq(5), fmpq(6), fmpq(9, 10), 2),
    (fmpq(6), fmpq(7), fmpq(19, 25), 2),
    (fmpq(7), fmpq(8), fmpq(33, 50), 2),
    (fmpq(8), fmpq(8803377, 1000000), fmpq(83, 100), 3),
]
WITNESS = (fmpq(88033777, 10000000), fmpq(107538996, 10000000))   # (alpha, x) = (8.8033777, 10.7538996)


def ball(lo, hi):
    """Ball containing the closed interval [lo, hi] (fmpq or arb endpoints)."""
    a = lo if isinstance(lo, arb) else arb(lo)
    b = hi if isinstance(hi, arb) else arb(hi)
    return a.union(b)


class Dual:
    """Value with enclosures of the partial derivatives d/dalpha and d/dx."""
    __slots__ = ('v', 'a', 'x')

    def __init__(self, v, a=None, x=None):
        self.v = v
        self.a = a if a is not None else arb(0)
        self.x = x if x is not None else arb(0)

    def __add__(self, o):
        o = lift(o)
        return Dual(self.v + o.v, self.a + o.a, self.x + o.x)

    __radd__ = __add__

    def __sub__(self, o):
        o = lift(o)
        return Dual(self.v - o.v, self.a - o.a, self.x - o.x)

    def __rsub__(self, o):
        return lift(o) - self

    def __mul__(self, o):
        o = lift(o)
        return Dual(self.v * o.v, self.a * o.v + self.v * o.a, self.x * o.v + self.v * o.x)

    __rmul__ = __mul__

    def __neg__(self):
        return Dual(-self.v, -self.a, -self.x)


def lift(o):
    if isinstance(o, Dual):
        return o
    return Dual(o if isinstance(o, arb) else arb(o))


def dexp(d):
    e = d.v.exp()
    return Dual(e, e * d.a, e * d.x)


def dsin(d):
    c = d.v.cos()
    return Dual(d.v.sin(), c * d.a, c * d.x)


def dcos(d):
    s = d.v.sin()
    return Dual(d.v.cos(), -s * d.a, -s * d.x)


def dinv(d):
    i = 1 / d.v
    return Dual(i, -(i * i) * d.a, -(i * i) * d.x)


def T_dual(A, X, K):
    """e^{-x} S(alpha, x) with derivative enclosures, for alpha in A, x in X."""
    al = Dual(A, arb(1), arb(0))
    xx = Dual(X, arb(0), arb(1))
    th = dinv(al) * (2 * PI)
    tot = Dual(arb(0))
    for k in range(1, K + 1):
        u = th * k
        c = dcos(u)
        s = dsin(u)
        e = dexp(-(xx * (1 - c)))
        p = xx * s
        tot = tot + e * (dcos(p) - u * xx * dsin(p + u)) * 2
    return tot


def m2_lower(A, phi):
    """Lower bound of m^2 = min_r |r^alpha e^{i alpha phi} - 1|^2: 1 if cos(alpha phi) <= 0, else sin^2(alpha phi)."""
    g = A * phi
    if g.cos().upper() <= 0:
        return arb(1)
    m2 = min(arb(1), (g.sin() ** 2).lower())
    assert m2 > 0
    return m2


def B_upper(A, X, phi):
    """Upper bound of alpha^2 e^{-x} B(alpha, x) over the box."""
    m2 = m2_lower(A, phi)
    lam = -X * phi.cos()
    assert lam > 0
    lnlam = lam.log()
    p1 = A.gamma() * (-(A * lnlam)).exp() * (phi + max((A.digamma() - lnlam).upper(), arb(0)))
    l1 = (1 / A ** 2).upper()
    l2 = ((A - 1).gamma() / (E * (lam ** (A - 1)))).upper()
    B = (p1 + 2 * min(l1, l2)) / (PI.lower() * m2)
    return ((A ** 2) * (-X).exp() * B).upper()


def check_range_geometry(lo, hi, phi_frac, K):
    A = ball(lo, hi)
    phi = PI * arb(phi_frac)
    th = 2 * PI / A
    assert phi > PI / 2 and phi < PI
    for k in range(1, K + 1):
        assert k * th < phi, ('pole should be inside', k)
    assert (K + 1) * th > phi, 'next pole must be outside the sector'
    n = 400
    for i in range(n):
        Ai = ball(lo + (hi - lo) * i / n, lo + (hi - lo) * (i + 1) / n)
        assert m2_lower(Ai, phi) > arb(1) / 2, 'rays must stay well away from the poles'
    return phi


def Q_lower(alo, ahi, xlo, xhi, phi, K):
    A = ball(alo, ahi)
    X = ball(xlo, xhi)
    tb = T_dual(A, X, K)
    naive = tb.v.lower()
    ac = arb((alo + ahi) / 2)
    xc = (xlo + xhi) / 2
    ra = arb((ahi - alo) / 2)
    rx = (xhi - xlo) / 2
    tc = T_dual(ac, xc, K).v
    centered = (tc - tb.a.abs_upper() * ra - tb.x.abs_upper() * rx).lower()
    t_low = max(naive, centered)
    q = (1 + t_low - B_upper(A, X, phi)).lower()
    return q, tb.a.abs_upper() * ra, tb.x.abs_upper() * rx


def certify_middle():
    stats = dict(accepted=0, evals=0, maxdepth=0, minq=arb(10), minq_box=None)
    for lo, hi, phi_frac, K in RANGES:
        phi = check_range_geometry(lo, hi, phi_frac, K)
        n_slabs = int(((hi - lo) * 20).floor()) or 1           # slabs of width 1/20
        for i in range(n_slabs):
            alo = lo + (hi - lo) * i / n_slabs
            ahi = lo + (hi - lo) * (i + 1) / n_slabs
            xlo = (1 + arb(alo)).digamma().exp().lower()       # <= e^{psi(1+alpha)} on the slab
            stack = [(alo, ahi, xlo, X_TOP, 0)]
            while stack:
                a0, a1, x0, x1, depth = stack.pop()
                if depth > 40:
                    print('FAIL: depth guard at', a0, a1, x0, x1)
                    return None
                stats['evals'] += 1
                q, ea, ex = Q_lower(a0, a1, x0, x1, phi, K)
                if q > 0:
                    stats['accepted'] += 1
                    stats['maxdepth'] = max(stats['maxdepth'], depth)
                    if q < stats['minq']:
                        stats['minq'] = q
                        stats['minq_box'] = (a0, a1, x0, x1)
                    continue
                if ea > ex:
                    am = (a0 + a1) / 2
                    stack += [(a0, am, x0, x1, depth + 1), (am, a1, x0, x1, depth + 1)]
                else:
                    xm = (x0 + x1) / 2
                    stack += [(a0, a1, x0, xm, depth + 1), (a0, a1, xm, x1, depth + 1)]
    return stats


def certify_tail():
    worst = arb(0)
    for lo, hi, phi_frac, K in RANGES:
        phi = check_range_geometry(lo, hi, phi_frac, K)
        n = int(((hi - lo) * 200).floor()) or 1
        for i in range(n):
            A = ball(lo + (hi - lo) * i / n, lo + (hi - lo) * (i + 1) / n)
            th = 2 * PI / A
            g = arb(0)
            for k in range(1, K + 1):
                u = k * th
                g += 2 * (-(X_TOP * (1 - u.cos()))).exp() * (1 + u * X_TOP)
            G = (g.upper() + B_upper(A, X_TOP, phi)).upper()
            assert G < 1, ('tail bound fails', i)
            worst = max(worst, G)
    return worst


def numeric_representation_check():
    """Non-rigorous sanity check of the residue + ray representation against the series (mpmath, 40 digits)."""
    import mpmath as mp
    mp.mp.dps = 40

    def D_series(a, x):
        a, x = mp.mpf(a), mp.mpf(x)
        lx, s, n = mp.log(x), mp.mpf(0), 1
        while True:
            an = a * n
            t = mp.exp(an * lx - mp.loggamma(1 + an))
            s += n * (lx - mp.digamma(1 + an)) * t
            if an > x + 60 and t < mp.mpf(10) ** -50:
                return s
            n += 1

    def D_rep(a, x, phi):
        a, x, phi = mp.mpf(a), mp.mpf(x), mp.mpf(phi)
        th, S, k = 2 * mp.pi / a, mp.mpf(0), 1
        while k * th < phi:
            u = k * th
            S += 2 * mp.exp(x * mp.cos(u)) * (mp.cos(x * mp.sin(u)) - u * x * mp.sin(x * mp.sin(u) + u))
            k += 1
        e, eap = mp.exp(1j * phi), mp.exp(1j * a * phi)
        f = lambda r: mp.exp(x * r * e) * r ** (a - 1) * eap * (mp.log(r) + 1j * phi) / (r ** a * eap - 1) ** 2
        J = mp.im(mp.quad(f, [0, 0.5, 1, 2, 5, 20, mp.inf])) / mp.pi
        return -(mp.exp(x) + S) / a ** 2 - J

    worst = mp.mpf(0)
    for a, pf in [(4.5, 0.65), (5.5, 0.9), (6.0, 0.76), (6.5, 0.76), (7.5, 0.66), (8.0, 0.83), (8.5, 0.83), (8.8, 0.83)]:
        for x in [1.5, 6.0, 10.75, 14.0]:
            ds = D_series(a, x)
            worst = max(worst, abs(ds - D_rep(a, x, pf * mp.pi)) / abs(ds))
    print(f'representation check (32 cases, alpha in [4.5, 8.8]): max relative difference = {mp.nstr(worst, 3)}')


def witness(alpha, x, N=40):
    """Rigorous enclosure of D(alpha, x) from the series: partial sum to N plus tail bound."""
    A = arb(alpha)
    X = arb(x)
    lx = X.log()
    s = arb(0)
    for n in range(1, N + 1):
        an = A * n
        s += n * (lx - (1 + an).digamma()) * (an * lx - (1 + an).lgamma()).exp()
    # tail: |term_n| <= n (log(1+an) + |log x|) (e x / (a n))^(a n)  (Stirling: Gamma(1+y) >= (y/e)^y)
    bnd = []
    for n in (N + 1, N + 2):
        an = A * n
        bnd.append(n * ((1 + an).log() + lx.abs_upper()) * ((E * X / an) ** an))
    ratio = (bnd[1] / bnd[0]).upper()
    assert ratio < arb(1) / 2
    tail = (2 * bnd[0]).upper()
    return s, tail


if __name__ == '__main__':
    # optional:  --numeric (representation vs series, ~2 min)  --amax p/q  (top of the certified alpha range)   --witness p/q,r/s  (alpha, x) of the failure witness
    args = sys.argv[1:]
    if '--amax' in args:
        p, q = args[args.index('--amax') + 1].split('/')
        lo, _, ph, K = RANGES[-1]
        RANGES[-1] = (lo, fmpq(int(p), int(q)), ph, K)
    if '--witness' in args:
        wa, wx = args[args.index('--witness') + 1].split(',')
        WITNESS = tuple(fmpq(*map(int, w.split('/'))) for w in (wa, wx))
    t0 = time.time()
    tail_worst = certify_tail()
    print(f'certified alpha range: [4, {float(RANGES[-1][1])}]')
    print(f'x >= 16: analytic bound G < 1 on the whole range; max certified G(alpha,16) = {tail_worst.str(5, radius=False)}')
    st = certify_middle()
    if st is None:
        print('MIDDLE REGION FAILED')
        sys.exit(1)
    b = st['minq_box']
    print(f"middle region e^psi(1+a) <= x <= 16: accepted boxes = {st['accepted']}, evaluations = {st['evals']}, "
          f"max depth = {st['maxdepth']}")
    print(f"   smallest certified lower bound of Q = {st['minq'].str(5, radius=False)} on alpha in "
          f"[{float(b[0]):.6f}, {float(b[1]):.6f}], x in [{float(b[2].mid()):.5f}, {float(b[3].mid()):.5f}]")
    s, tail = witness(*WITNESS)
    print(f'witness D({float(WITNESS[0])}, {float(WITNESS[1])}): partial sum = {s.str(10)}, tail <= {tail.str(3, radius=False)}; '
          f'certified positive: {(s - tail) > 0}')
    print(f'elapsed {time.time() - t0:.1f} s')
    if '--numeric' in args:
        numeric_representation_check()
\end{lstlisting}



\subsection{No recovery (\texttt{verify\_no\_recovery.py})}
\begin{lstlisting}
from flint import arb, fmpq, ctx
ctx.prec=128
E=arb(1).exp()
c=arb(fmpq(5,4))
b2=5*E/16
b3=25*E**2/432
delta=c.log()-arb(fmpq(1,18))-(arb(fmpq(1,108))).exp()*(arb(2).sqrt()*b2**9*((arb(fmpq(8,5))).log()+arb(fmpq(1,36)))+2*arb(3).sqrt()*b3**9*((arb(fmpq(12,5))).log()+arb(fmpq(1,54))))
assert delta>arb(fmpq(2459,1000000))
print('delta =',delta)
# D(alpha, x0) > 0 for all alpha in [8.8033777,9].
x=arb(fmpq(134423745,12500000)); lx=x.log()
N=20
j=N+1
# The discarded terms have total absolute value <= TAIL.
TAIL=26*j*j*arb(fmpq(33,8*j))**(8*j)
assert 4*arb(fmpq(33,8*(j+1)))**8<arb(fmpq(1,2))
stack=[(fmpq(88033777,10000000),fmpq(9),0)]
accepted=evals=maxdepth=0
minlb=arb(1000)
while stack:
 lo,hi,d=stack.pop(); assert d<=50
 A=arb(lo).union(arb(hi))
 D=arb(0)
 for n in range(1,N+1):
  an=n*A
  D += n*(an*lx-(1+an).lgamma()).exp()*(lx-(1+an).digamma())
 lb=(D-TAIL).lower(); evals+=1
 if lb>0:
  accepted+=1;maxdepth=max(maxdepth,d);minlb=min(minlb,lb)
 else:
  mid=(lo+hi)/2
  stack.extend([(lo,mid,d+1),(mid,hi,d+1)])
print('tail_bound =',TAIL)
print('accepted =',accepted,'evals =',evals,'maxdepth =',maxdepth,'min lower bound =',minlb)

\end{lstlisting}

\subsection{Second parameter two (\texttt{verify\_beta2\_range.py})}
\begin{lstlisting}
#!/usr/bin/env python3
"""Arb certificate for -d/dalpha E_{alpha,2}(x^alpha)>0 on alpha in [4,9], x>0.
Small x follows analytically from x<=exp(psi(alpha+2)).
The script covers the middle region and x>=16, using residue/ray bounds
for the inverse Laplace transform of x H_{alpha,2}(x).
Requires python-flint; tested with version0.9.0, precision128bits.
"""
import sys
import time
from flint import arb, fmpq, ctx

ctx.prec = 128
PI = arb.pi()
E = arb(1).exp()
X_TOP = arb(16)

# (alpha_lo, alpha_hi, phi / pi, K = number of residue pairs inside |arg w| < phi)
RANGES = [
    (fmpq(4), fmpq(5), fmpq(13, 20), 1),
    (fmpq(5), fmpq(6), fmpq(9, 10), 2),
    (fmpq(6), fmpq(7), fmpq(19, 25), 2),
    (fmpq(7), fmpq(8), fmpq(33, 50), 2),
    (fmpq(8), fmpq(9), fmpq(83, 100), 3),
]


def ball(lo, hi):
    """Ball containing the closed interval [lo, hi] (fmpq or arb endpoints)."""
    a = lo if isinstance(lo, arb) else arb(lo)
    b = hi if isinstance(hi, arb) else arb(hi)
    return a.union(b)


class Dual:
    """Value with enclosures of the partial derivatives d/dalpha and d/dx."""
    __slots__ = ('v', 'a', 'x')

    def __init__(self, v, a=None, x=None):
        self.v = v
        self.a = a if a is not None else arb(0)
        self.x = x if x is not None else arb(0)

    def __add__(self, o):
        o = lift(o)
        return Dual(self.v + o.v, self.a + o.a, self.x + o.x)

    __radd__ = __add__

    def __sub__(self, o):
        o = lift(o)
        return Dual(self.v - o.v, self.a - o.a, self.x - o.x)

    def __rsub__(self, o):
        return lift(o) - self

    def __mul__(self, o):
        o = lift(o)
        return Dual(self.v * o.v, self.a * o.v + self.v * o.a, self.x * o.v + self.v * o.x)

    __rmul__ = __mul__

    def __neg__(self):
        return Dual(-self.v, -self.a, -self.x)


def lift(o):
    if isinstance(o, Dual):
        return o
    return Dual(o if isinstance(o, arb) else arb(o))


def dexp(d):
    e = d.v.exp()
    return Dual(e, e * d.a, e * d.x)


def dsin(d):
    c = d.v.cos()
    return Dual(d.v.sin(), c * d.a, c * d.x)


def dcos(d):
    s = d.v.sin()
    return Dual(d.v.cos(), -s * d.a, -s * d.x)


def dinv(d):
    i = 1 / d.v
    return Dual(i, -(i * i) * d.a, -(i * i) * d.x)


def T_dual(A, X, K):
    """e^{-x} S_2(alpha, x) with derivative enclosures, for alpha in A, x in X."""
    al = Dual(A, arb(1), arb(0))
    xx = Dual(X, arb(0), arb(1))
    th = dinv(al) * (2 * PI)
    tot = Dual(arb(0))
    for k in range(1, K + 1):
        u = th * k
        c = dcos(u)
        s = dsin(u)
        e = dexp(-(xx * (1 - c)))
        p = xx * s
        tot = tot + e * (dcos(p - u) + u * dsin(p - u) - u * xx * dsin(p)) * 2
    return tot


def m2_lower(A, phi):
    """Lower bound of m^2 = min_r |r^alpha e^{i alpha phi} - 1|^2: 1 if cos(alpha phi) <= 0, else sin^2(alpha phi)."""
    g = A * phi
    if g.cos().upper() <= 0:
        return arb(1)
    m2 = min(arb(1), (g.sin() ** 2).lower())
    assert m2 > 0
    return m2


def B_upper(A, X, phi):
    """Upper bound of alpha^2 e^{-x} B(alpha, x) over the box."""
    m2 = m2_lower(A, phi)
    lam = -X * phi.cos()
    assert lam > 0
    lnlam = lam.log()
    p1 = (A - 1).gamma() * (-((A - 1) * lnlam)).exp() * (phi + max(((A - 1).digamma() - lnlam).upper(), arb(0)))
    l1 = (1 / (A - 1) ** 2).upper()
    l2 = ((A - 2).gamma() / (E * (lam ** (A - 2)))).upper()
    B = (p1 + 2 * min(l1, l2)) / (PI.lower() * m2)
    return ((A ** 2) * (-X).exp() * B).upper()


def check_range_geometry(lo, hi, phi_frac, K):
    A = ball(lo, hi)
    phi = PI * arb(phi_frac)
    th = 2 * PI / A
    assert phi > PI / 2 and phi < PI
    for k in range(1, K + 1):
        assert k * th < phi, ('pole should be inside', k)
    assert (K + 1) * th > phi, 'next pole must be outside the sector'
    n = 400
    for i in range(n):
        Ai = ball(lo + (hi - lo) * i / n, lo + (hi - lo) * (i + 1) / n)
        assert m2_lower(Ai, phi) > arb(1) / 2, 'rays must stay well away from the poles'
    return phi


def Q_lower(alo, ahi, xlo, xhi, phi, K):
    A = ball(alo, ahi)
    X = ball(xlo, xhi)
    tb = T_dual(A, X, K)
    naive = tb.v.lower()
    ac = arb((alo + ahi) / 2)
    xc = (xlo + xhi) / 2
    ra = arb((ahi - alo) / 2)
    rx = (xhi - xlo) / 2
    tc = T_dual(ac, xc, K).v
    centered = (tc - tb.a.abs_upper() * ra - tb.x.abs_upper() * rx).lower()
    t_low = max(naive, centered)
    q = (1 + t_low - B_upper(A, X, phi)).lower()
    return q, tb.a.abs_upper() * ra, tb.x.abs_upper() * rx


def certify_middle():
    stats = dict(accepted=0, evals=0, maxdepth=0, minq=arb(10), minq_box=None)
    for lo, hi, phi_frac, K in RANGES:
        phi = check_range_geometry(lo, hi, phi_frac, K)
        n_slabs = int(((hi - lo) * 20).floor()) or 1           # exact rational slabs of width 1/20
        for i in range(n_slabs):
            alo = lo + (hi - lo) * i / n_slabs
            ahi = lo + (hi - lo) * (i + 1) / n_slabs
            xlo = (2 + arb(alo)).digamma().exp().lower()       # <= e^{psi(2+alpha)} on the slab
            stack = [(alo, ahi, xlo, X_TOP, 0)]
            while stack:
                a0, a1, x0, x1, depth = stack.pop()
                if depth > 40:
                    print('FAIL: depth guard at', a0, a1, x0, x1)
                    return None
                stats['evals'] += 1
                q, ea, ex = Q_lower(a0, a1, x0, x1, phi, K)
                if q > 0:
                    stats['accepted'] += 1
                    stats['maxdepth'] = max(stats['maxdepth'], depth)
                    if q < stats['minq']:
                        stats['minq'] = q
                        stats['minq_box'] = (a0, a1, x0, x1)
                    continue
                if ea > ex:
                    am = (a0 + a1) / 2
                    stack += [(a0, am, x0, x1, depth + 1), (am, a1, x0, x1, depth + 1)]
                else:
                    xm = (x0 + x1) / 2
                    stack += [(a0, a1, x0, xm, depth + 1), (a0, a1, xm, x1, depth + 1)]
    return stats


def certify_tail():
    worst = arb(0)
    for lo, hi, phi_frac, K in RANGES:
        phi = check_range_geometry(lo, hi, phi_frac, K)
        n = int(((hi - lo) * 200).floor()) or 1
        for i in range(n):
            A = ball(lo + (hi - lo) * i / n, lo + (hi - lo) * (i + 1) / n)
            th = 2 * PI / A
            g = arb(0)
            for k in range(1, K + 1):
                u = k * th
                g += 2 * (-(X_TOP * (1 - u.cos()))).exp() * (1 + u + u * X_TOP)
            G = (g.upper() + B_upper(A, X_TOP, phi)).upper()
            assert G < 1, ('tail bound fails', i)
            worst = max(worst, G)
    return worst


if __name__ == '__main__':
    t0 = time.time()
    print('BETA 2 CERTIFICATE: alpha in [4,9]')
    tail = certify_tail()
    assert tail < arb('0.610')
    print('Maximum tail-region G bound:', tail.str(25))
    st = certify_middle()
    assert st is not None
    assert st['minq'] > 0
    print('Middle-region certificate:', st)
    print('PASS; elapsed seconds:',time.time()-t0)

\end{lstlisting}

\end{document}
